% ======================================================================
\section{Finite terms}\label{sec:finite}
% ======================================================================
% The equations of the rows (files rows/*.tex) are generated mechanically by scripts/gen_rows.py from
% Sec. 8 of nlo_review/final_explicit/Final_terms_explicit_2_fixed.tex; do not edit them by hand.

After the rapidity logarithms, the poles, and the DGLAP and running-coupling terms have been taken out, what remains of each row is finite. The finite terms are, row by row, the limit $\Lambda\to0$, $\vee\to\infty$ of
\begin{align}\label{eq:fin_def}
&\big[\text{row with the cut-offs}\big]\notag\\&\quad-\log\frac{\vee}{k^+}\times\big[\text{its term in the BFKL list}\big]-\log\frac{k^+}{\Lambda}\times\big[\text{its term in the JIMWLK list}\big]\,,
\end{align}
where the lists are those whose sums are Eqs.~\eqref{eq:BFKLrho} and~\eqref{eq:JIMWLKrho}. The collinear poles are treated as in Section~\ref{sec:dglap_rc}, and the DGLAP and running-coupling terms of Row~II.IV are those of Section~\ref{sec:dg_final}. Finite terms arise in Rows I.II, I.III, I.V, II.I, II.IV, III.II, III.III, III.V and III.VI. The rows with four charges only (I.I, II.II, II.III, III.IV) contribute to the evolution terms only, and the finite parts of Rows~I.IV and~III.I cancel identically.

In this section we give the finite terms Group by Group and Row by Row, in two forms. In the first form the cancellations that are exact have been carried out and every product of three charges has been written as its totally symmetrised product plus commutator terms, which contain two charges; this form holds for an arbitrary projectile and is given in Appendix~\ref{app:finite}. In the second form the result is averaged over the projectile and written in terms of colour-singlet combinations of Wilson lines, Section~\ref{sec:fin_rows}. In both forms the Wilson lines appear, as at leading order, through differences $U-U$, and the longitudinal integrals are written with plus prescriptions in a momentum fraction $\xi$. We begin with a summary of the complete result, then state the conventions, the treatment of the ordered charges and of the average over the projectile, and the conditions under which the terms are integrable, and finally list the terms.

\subsection{Summary of the result}\label{sec:result_summary}

The NLO single inclusive gluon spectrum, Eq.~\eqref{eq:schematic}, is
\begin{align}\label{eq:result}
\frac{d^3N}{d^3k}&=\Big[\frac{d^3N_{\rm LO}}{d^3k}+\frac{d^3N_{\rm RC}}{d^3k}\Big]+\log\frac{\vee}{k^+}\,\frac{d^3N_{\rm BFKL}}{d^3k}+\log\frac{k^+}{\Lambda}\,\frac{d^3N_{\rm JIMWLK}}{d^3k}\notag\\
&\quad+\frac{d^3N_{\rm DGLAP}}{d^3k}+\sum_{\rm rows}\Big[\frac{d^3N^{\rm row}_{\rm sym}}{d^3k}+\frac{d^3N^{\rm row}_{\rm con}}{d^3k}+\frac{d^3N^{\rm row}_{\rm quad}}{d^3k}\Big]\,,
\end{align}
in pure Yang--Mills theory ($N_f=0$), in the $\MS$ scheme with conventional dimensional regularisation (Section~\ref{sec:dg_poles}), with the $\MS$ number-density fragmentation function (Section~\ref{sec:dg_msbar}), at leading power in $k^+/\vee$ and $\Lambda/k^+$ (eikonal approximation, Section~\ref{sec:dg_poles}). The finite terms of a row consist of up to three parts: the symmetrised three-$\rho$ part (``sym''), the two-$\rho$ contact terms that the symmetrisation of the three charges leaves behind (``con''), and the genuine two-$\rho$ terms (``quad''). The pieces, with pointers to all equations, are collected in Table~\ref{tab:result}.
\begin{table}[t]
\centering
\renewcommand{\arraystretch}{1.18}
\footnotesize\setlength{\tabcolsep}{3pt}
\begin{tabular}{@{}>{\raggedright\arraybackslash}p{0.17\textwidth}>{\raggedright\arraybackslash}p{0.27\textwidth}>{\raggedright\arraybackslash}p{0.27\textwidth}>{\raggedright\arraybackslash}p{0.24\textwidth}@{}}
\toprule
term & ordered charges & averaged, colour singlet & content\\
\midrule
LO $+$ RC & \eqref{eq:LO}, \eqref{eq:LORC}; local form \eqref{eq:dg_local} & \eqref{eq:LOsinglet}, \eqref{eq:sLOrepl} & LO with $\alpha_s(c_0^2/\v r^2)$, resp.\ $\alpha_s(\v k^2)$; separate $\mu_R$, $\mu_F$: \eqref{eq:dg_muRmuF}\\
BFKL & \eqref{eq:BFKLrho} & \eqref{eq:BFKLsinglet} & one LO BFKL step of the projectile correlator, \eqref{eq:BFKLasLO}\\
JIMWLK & \eqref{eq:JIMWLKrho} & \eqref{eq:JIMWLKsinglet} & $-H_{\rm JIMWLK}$ on the LO target operator, \eqref{eq:JIMWLKasLO}\\
DGLAP & \eqref{eq:DGLAP}; local form \eqref{eq:dg_local}, \eqref{eq:dg_R} & \eqref{eq:sLOrepl} & final-state $\MS$ fragmentation with $\widetilde{\mathcal P}_{gg}$\\
\midrule
Rows I.I, II.II, II.III, III.IV & --- & --- & evolution terms only\\
Rows I.IV $+$ III.I & --- & --- & finite parts cancel identically\\
Row I.II & \eqref{eq:H-I.II-3}, \eqref{eq:H-I.II-2}, \eqref{eq:H-I.II-2o} & \eqref{eq:S-I.II-3}, \eqref{eq:S-I.II-2}, \eqref{eq:S-I.II-2o} & vertex correction, $\mf A^{(3)}$\\
Row I.III & \eqref{eq:H-I.III-3}, \eqref{eq:H-I.III-2} & \eqref{eq:S-I.III-3}, \eqref{eq:S-I.III-2} & rescattering, $\mf B_3$\\
Row I.V & \eqref{eq:H-I.V-3}, \eqref{eq:H-I.V-2}, \eqref{eq:H-I.V-2o} & \eqref{eq:S-I.V-3}, \eqref{eq:S-I.V-2}, \eqref{eq:S-I.V-2o} & recombination, $\mf C_1^\dagger\mf B_2$\\
Row II.I & \eqref{eq:H-II.I-3}, \eqref{eq:H-II.I-2}, \eqref{eq:H-II.I-2a}--\eqref{eq:H-II.I-2e} & \eqref{eq:S-II.I-3}, \eqref{eq:S-II.I-2}, \eqref{eq:S-II.I-2a}--\eqref{eq:S-II.I-2e} & $|\mathcal A_B|^2$\\
Row II.IV & \eqref{eq:H-II.IVrem} & \eqref{eq:S-II.IVrem} & $|\mathcal A_{CA}|^2$: tensor term and remnant\\
Row III.II & \eqref{eq:H-III.II-3}, \eqref{eq:H-III.II-2} & \eqref{eq:S-III.II-3}, \eqref{eq:S-III.II-2} & interference\\
Row III.III & \eqref{eq:H-III.III-3}, \eqref{eq:H-III.III-2}, \eqref{eq:H-III.III-2o} & \eqref{eq:S-III.III-3}, \eqref{eq:S-III.III-2}, \eqref{eq:S-III.III-2o} & $-2{\rm Re}\,\mathcal A_{CA}^\dagger\mathcal A_B$\\
Row III.V & \eqref{eq:H-III.V-3}, \eqref{eq:H-III.V-2} & \eqref{eq:S-III.V-3}, \eqref{eq:S-III.V-2} & interference\\
Row III.VI & \eqref{eq:H-III.VI-3}, \eqref{eq:H-III.VI-2} & \eqref{eq:S-III.VI-3}, \eqref{eq:S-III.VI-2} & interference\\
\bottomrule
\end{tabular}
\caption{The complete NLO result, Eq.~\eqref{eq:result}: where each term is given, for ordered charges of an arbitrary projectile and averaged over the projectile in colour-singlet form. For each row the equations are listed in the order sym, con, quad.}\label{tab:result}
\end{table}
For each row the equations of Table~\ref{tab:result} are listed in the order sym, con, quad. The first two parts are present in the rows that are cubic in the charges (we also call them the first, or cubic, part of a row, and the genuine two-$\rho$ terms its second part); Rows~I.II, I.V, II.I and~III.III have in addition, and Row~II.IV has only, genuine two-$\rho$ terms. After the average over the projectile (Section~\ref{sec:fin_singlet}) the symmetrised part is proportional to the symmetrised three-point function $\sigma$ of the projectile, Eq.~\eqref{eq:A2d}, and is odd in $\v k$, while the contact and the genuine two-$\rho$ terms are proportional to the two-point function $n$ and are even in $\v k$. The prescriptions that define the finite terms are part of the result: the vanishing energy denominators of Row~I.III are principal values, and certain terms, whose large-distance tails cancel only in combination, have to be integrated together (Section~\ref{sec:fin_together}). The finite terms are given in coordinate space; no numerical evaluation of them, and no plot, is given in this paper.

\paragraph{Assembling the spectrum.}\label{sec:howtoread} To obtain the spectrum for a given projectile and target:
\begin{enumerate}
\item Choose the form: ordered charges (Appendix~\ref{app:finite}, any projectile, including the one-$\rho$ terms of Eqs.~\eqref{eq:symord}--\eqref{eq:symr21} if $\langle\rho\rangle\neq0$) or averaged (Section~\ref{sec:fin_rows}, with the assumptions (A1), (A2) of Section~\ref{sec:fin_singlet}).
\item Take LO $+$ RC, the BFKL and JIMWLK terms and the DGLAP term from the first block of Table~\ref{tab:result}. The DGLAP and RC terms come in a coordinate-space and a local form, Eq.~\eqref{eq:dg_equiv}; use one of them consistently. The remnant is counted once: either as the second term of Eq.~\eqref{eq:S-II.IVrem} with Eq.~\eqref{eq:DGLAP}, or inside Eq.~\eqref{eq:DGLAPdoubleplus}, or inside $N_{\mathcal R}$, Eq.~\eqref{eq:dg_R}.
\item Add the finite terms of the nine rows of the second block of Table~\ref{tab:result}, each with its sym, con and quad parts where present.
\item Respect the groupings and orders of integration of Section~\ref{sec:fin_together}: the four pairs A--D under one $\v z$ integral; the $\v w$ integral before the $\v z$ integral in Row~I.III; Rows~II.IV, II.I (quad) and III.III (quad) under one integral over the position of the pair; the instantaneous terms with the $\v y'$ integral before the $\xi$ integral; the collinear point of Row~I.V with a rotation-invariant cut-off.
\item For a comparison with data or with a leading-logarithmic calculation, absorb the two rapidity logarithms into the evolution of the projectile and target averages (end of Section~\ref{sec:logsummary}), and convolute with the $\MS$ fragmentation function at the scale $\mu_F$ of the DGLAP term.
\end{enumerate}

\paragraph{Where the terms are.} Terms even in $\v k$: LO, RC, BFKL, JIMWLK, DGLAP, all quad and all con terms. Terms odd in $\v k$: the sym terms, and only they. Terms with a scale $\mu$: RC and DGLAP only. Terms with $\vee$ or $\Lambda$: BFKL and JIMWLK only. Terms that depend on the principal-value prescription: Row~I.III only.

\subsection{Conventions}\label{sec:fin_conv}

\paragraph{Phase and notation.} All terms are written with the phase $e^{-i\v k\cdot(\text{unprimed}-\text{primed})}$, with the conjugated factor at the primed points and to the left, e.g.\ $e^{-i\v k\cdot(\v w-\v w')}$. In the rows in which the measured gluon is part of a two-gluon state, the phase contains a recoil factor: $e^{-i\v k\cdot(\v z-\v w)\xi}$ in Row~I.III, $e^{-i\v k\cdot(\v z-\v w')}e^{-i\v k\cdot(\v x-\v z)\xi}$ in Row~I.V, $e^{-i\v k\cdot(\v w-\v y')}e^{-i\v k\cdot(\v z-\v y')\xi}$ in Rows~III.III, III.V and~III.VI, and $e^{-i\v k\cdot(\v y-\v y')/\xi}$ in Row~II.IV and in the DGLAP term. The convention matters: a change of convention in a row amounts to $\v k\to-\v k$ there, which leaves the part of a term that is even in $\v k$ unchanged and flips the sign of its odd part, i.e.\ of the symmetrised three-$\rho$ part (Section~\ref{sec:fin_singlet}); mixing conventions between rows would give wrong relative signs of the three-$\rho$ terms. The complex conjugate is implied for Groups~I and~III and is written as ``$+{\rm c.c.}$''; Group~II has none. The charges $\rho$ are operators in the order written. In the formulas of this section and of Appendix~\ref{app:finite} transverse vectors are not printed in bold, $\int_x\equiv\int d^2x$, $\Kv(x)=x/x^2$ and $K(x,y,z)=\Kv(x-z)\cdot\Kv(y-z)$ as in Eq.~\eqref{eq:kernels}, and $\{\dots\}_{\rm sym}$ denotes the average of a product of charges over all orderings. The adjoint Wilson lines $U^{ab}$ are real, $U^\dagger=U^T$, and
\begin{align}\label{eq:TD}
(T^e)_{bc}=-if^{ebc}\,,\qquad (D^e)_{bc}=d^{ebc}\,,\qquad{\rm Tr}\,T^aT^b=N_c\,\delta^{ab}\,.
\end{align}

\paragraph{Longitudinal momentum fraction.} The longitudinal integrals are written in terms of a momentum fraction $\xi$, with two definitions. In all rows except Row~II.IV,
\begin{align}\label{eq:xi_def}
\xi=\frac{p^+}{k^+}\,,
\end{align}
where $p^+$ is the longitudinal momentum of the gluon that is integrated over: $0\le\xi\le1$ in Group~I, where this gluon is part of the measured one, and $0\le\xi<\infty$ in Group~III, where the upper limit $\vee-k^+$ has been sent to infinity. In Row~II.IV, as in the DGLAP term of Section~\ref{sec:dg_final},
\begin{align}\label{eq:xi_defIV}
\xi=\frac{k^+}{p^++k^+}
\end{align}
is the fraction of the longitudinal momentum of the parent gluon that is carried by the measured gluon; then $p^+=k^+\frac{1-\xi}{\xi}$, $dp^+=-k^+d\xi/\xi^2$, and the range $0\le p^+\le\vee-k^+$ is $k^+/\vee\le\xi\le1$. Every row has the overall factor $1/k^+$ of the leading-order spectrum.

\paragraph{Plus prescriptions.} The plus prescriptions act on the whole $\xi$-dependent function $f$ of the term in which they stand, i.e.\ on the recoil phase and on the energy denominators:
\begin{align}\label{eq:plus_xi}
\int_0^1d\xi\,\Big[\frac{1}{\xi}\Big]_+f(\xi)&=\int_0^1\frac{d\xi}{\xi}\,\big[f(\xi)-f(0)\big]\,,\notag\\
\int_0^1d\xi\,\Big[\frac{1}{1-\xi}\Big]_+f(\xi)&=\int_0^1\frac{d\xi}{1-\xi}\,\big[f(\xi)-f(1)\big]\,,\notag\\
\int_0^\infty d\xi\,\Big[\frac{1}{\xi\,(1+\xi)}\Big]_+f(\xi)&=\int_0^\infty\frac{d\xi}{\xi\,(1+\xi)}\,\big[f(\xi)-f(0)\big]\,.
\end{align}
Where a plus prescription at $\xi\to0$ arises from the subtraction of Eq.~\eqref{eq:fin_def}, the subtracted term is the term of the row in the JIMWLK list.

\paragraph{Instantaneous terms of Group III.} The terms of the kernel $\mathfrak K^{(y')}_{\rm III}$ and the kernel $\mathfrak K^{(z)}_{\rm III}$ below that carry no factor depending on $\xi$ other than the phase are defined as distributions,
\begin{align}\label{eq:inst_xi}
\int_0^\infty d\xi\,e^{-i\v k\cdot(\v z-\v y')\,\xi}=\pi\,\delta(s)-i\,\mathcal P\frac1s\,,\qquad s=\v k\cdot(\v z-\v y')\,,
\end{align}
which act on the transverse kernels, with the two instantaneous terms of a row (those with $\mathfrak K^{(y')}_{\rm III}$ and with $\mathfrak K^{(z)}_{\rm III}$) under one $\v y'$ integral, done before the $\xi$ integral (Appendix~\ref{app:tails}).

\paragraph{Limits and their order.} The limits $\Lambda\to0$ and $\vee\to\infty$ are taken after the logarithmic (BFKL and JIMWLK) terms have been taken out, Eq.~\eqref{eq:fin_def}; terms of order $\Lambda/k^+$ and $k^+/\vee$ are dropped. The limit $\vee\to\infty$ is taken before $\epsilon\to0$; this order is part of the eikonal approximation (Section~\ref{sec:dg_poles}). Row~II.IV is written with its lower limit $\xi=k^+/\vee$; the limit $\vee\to\infty$ of this term exists once the integral over $\v y-\v y'$ has been done (see below Eq.~\eqref{eq:S-II.IVrem}).

\paragraph{Vanishing energy denominators.} The energy denominators of $\mf B_3$, which enter Row~I.III and its complex conjugate, vanish inside the region of integration. They are defined as principal values, Eq.~\eqref{eq:PV}; the finite terms of this paper refer to this prescription throughout. In the kernels of Row~I.III as printed the poles of the separate terms have cancelled, and all logarithms have positive arguments. An $i\varepsilon$ prescription differs by on-shell terms, which would add to Row~I.III a finite term that is not included here (Section~\ref{sec:nuance_pv}).

\subsection{Ordered charges: symmetrised products and contact terms}\label{sec:fin_sym}

The rows that are cubic in the charges contain ordered products $\rho\rho\rho$, $\rho\{\rho,\rho\}$ or $\{\rho,\rho\}\rho$; the order is the one in which the coefficients of $\bar\Omega^\dagger\Omega$ and of the conjugate amplitude act, with the conjugate amplitude to the left, and is not a choice. With the commutator of the charges, Eq.~\eqref{eq:rhoalgebra}, an ordered product of three charges is exactly the totally symmetrised product, plus single commutators, which are products of \emph{two} charges at coincident transverse positions, plus double commutators, which contain one charge. For any $c$-number colour tensor $C^{abc}$ multiplying the three charges (the Wilson lines are inside $C^{abc}$),
\begin{align}\label{eq:symord}
&C^{abc}\rho^{a}(p_{1})\rho^{b}(p_{2})\rho^{c}(p_{3})=C^{abc}\{\rho^{a}(p_{1})\rho^{b}(p_{2})\rho^{c}(p_{3})\}_{\rm sym}\notag\\
&\quad +\frac{i}{2}C^{abc}\Big[f^{abr}\delta^{(2)}(p_{1}-p_{2})\{\rho^{r}(p_{1})\rho^{c}(p_{3})\}_{\rm sym} +f^{bcr}\delta^{(2)}(p_{2}-p_{3})\{\rho^{a}(p_{1})\rho^{r}(p_{2})\}_{\rm sym}\notag\\
&\quad\quad +f^{acr}\delta^{(2)}(p_{1}-p_{3})\{\rho^{r}(p_{1})\rho^{b}(p_{2})\}_{\rm sym}\Big] -\frac{1}{6}C^{abc}\big[f^{abr}f^{rcs}\notag\\
&\quad\quad -f^{bcr}f^{ras}\big]\delta^{(2)}(p_{1}-p_{2})\, \delta^{(2)}(p_{1}-p_{3})\, \rho^{s}(p_{1})\,,
\end{align}

\begin{align}\label{eq:symr12}
&C^{abc}\rho^{a}(p_{1})\{\rho^{b}(p_{2}),\rho^{c}(p_{3})\}=2C^{abc}\{\rho^{a}(p_{1})\rho^{b}(p_{2})\rho^{c}(p_{3})\}_{\rm sym}\notag\\
&\quad +i\, C^{abc}\Big[f^{abr}\delta^{(2)}(p_{1}-p_{2})\{\rho^{r}(p_{1})\rho^{c}(p_{3})\}_{\rm sym} +f^{acr}\delta^{(2)}(p_{1}-p_{3})\{\rho^{r}(p_{1})\rho^{b}(p_{2})\}_{\rm sym}\Big]\notag\\
&\quad -\frac{1}{6}C^{abc}\big[f^{abr}f^{rcs} +f^{acr}f^{rbs}\big]\delta^{(2)}(p_{1}-p_{2})\, \delta^{(2)}(p_{1}-p_{3})\, \rho^{s}(p_{1})\,,
\end{align}

\begin{align}\label{eq:symr21}
&C^{abc}\{\rho^{a}(p_{1}),\rho^{b}(p_{2})\}\rho^{c}(p_{3})=2C^{abc}\{\rho^{a}(p_{1})\rho^{b}(p_{2})\rho^{c}(p_{3})\}_{\rm sym}\notag\\
&\quad +i\, C^{abc}\Big[f^{bcr}\delta^{(2)}(p_{2}-p_{3})\{\rho^{a}(p_{1})\rho^{r}(p_{2})\}_{\rm sym} +f^{acr}\delta^{(2)}(p_{1}-p_{3})\{\rho^{r}(p_{1})\rho^{b}(p_{2})\}_{\rm sym}\Big]\notag\\
&\quad +\frac{1}{6}C^{abc}\big[f^{bcr}f^{ras} +f^{acr}f^{rbs}\big]\delta^{(2)}(p_{1}-p_{2})\, \delta^{(2)}(p_{1}-p_{3})\, \rho^{s}(p_{1})\,.
\end{align}

Accordingly every cubic row is written as two equations: its \emph{symmetrised three-$\rho$ term}, and its \emph{two-$\rho$ terms from the symmetrisation} (contact terms), in which the $\delta$-functions put two of the three charges, and with them the corresponding Wilson lines of the colour tensor, at the same point. The one-$\rho$ terms, the last lines of Eqs.~\eqref{eq:symord}--\eqref{eq:symr21}, are not written in the rows; they vanish after the average over a projectile that is invariant under global colour rotations, $\langle\rho^s\rangle=0$, so that the averaged form of Section~\ref{sec:fin_rows} is complete, while as operator identities the rows of Appendix~\ref{app:finite} equal the ordered rows only together with them. The symmetrised term is the same operator for every ordering of the three charges; the ordering enters only through the contact terms. The contact terms are of the same order in the charges as the genuine two-$\rho$ terms, are present for any projectile, and are not removed by colour neutrality.

\subsection{Average over the projectile and colour-singlet form}\label{sec:fin_singlet}

\paragraph{Assumptions.} We assume
\begin{itemize}
\item[(A1)] \emph{Global colour invariance:} the projectile density matrix commutes with the total charge $Q^a=\int_x\rho^a(\v x)$. Then $\langle\rho^a(\v x)\rangle=0$, and the two-point function is colour diagonal,
\begin{align}\label{eq:A1n}
\big\langle\{\rho^a(\v x)\rho^b(\v y)\}_{\rm sym}\big\rangle=\frac{\delta^{ab}}{N_c^2-1}\,n(\v x,\v y)\,,\qquad n(\v x,\v y)\equiv\big\langle\rho^f(\v x)\rho^f(\v y)\big\rangle\,,
\end{align}
as in Eq.~\eqref{eq:singlet}, with $n(\v x,\v y)=n(\v y,\v x)$ real; the symmetry follows from $\langle[\rho^a(\v x),\rho^b(\v y)]\rangle=0$.
\item[(A2)] \emph{Symmetrised three-point function of $d$-type:}
\begin{align}\label{eq:A2d}
\big\langle\{\rho^a(\v x)\rho^b(\v y)\rho^c(\v z)\}_{\rm sym}\big\rangle&=d^{abc}\,\frac{N_c}{(N_c^2-4)(N_c^2-1)}\,\sigma(\v x,\v y,\v z)\,,\notag\\
\sigma(\v x,\v y,\v z)&\equiv\big\langle d^{efg}\{\rho^e(\v x)\rho^f(\v y)\rho^g(\v z)\}_{\rm sym}\big\rangle\,.
\end{align}
\end{itemize}
The normalisation follows from $d^{abc}d^{abc}=(N_c^2-4)(N_c^2-1)/N_c$; $\sigma$ is real and totally symmetric in its arguments. By (A1) the symmetrised three-point function is of the form $d^{abc}g_d+f^{abc}g_f$ with $g_d$ totally symmetric and $g_f$ totally antisymmetric in the three points; (A2) states that $g_f=0$ and gives $g_d$ the normalisation of Eq.~\eqref{eq:A2d}. The function $\sigma$ is odd under charge conjugation of the projectile: it is the three-point function of the cubic Casimir density, the odderon of the projectile. No model for $\sigma$ is assumed: $\sigma$ is a general real, totally symmetric function of three points, and all results below hold for any such $\sigma$. For $SU(2)$, $d^{abc}=0$ and the terms with $\sigma$ are absent. Nothing else is assumed; in particular no assumption on the four-point function is needed, since the rows with four charges contribute only to the evolution terms.

\paragraph{Quadratic and cubic terms.} For a genuine two-$\rho$ term, $\langle C^{ab}\rho^a(p)\rho^b(q)\rangle=C^{aa}\,n(p,q)/(N_c^2-1)$. For a cubic term, Eqs.~\eqref{eq:symord}--\eqref{eq:symr21} with (A1) and (A2) give, with $\Phi_{ij}\equiv C^{abc}f^{abc}$ evaluated with the Wilson lines at $p_i=p_j$,
\begin{align}\label{eq:cubicavg}
&\langle C^{abc}\rho^a(p_1)\rho^b(p_2)\rho^c(p_3)\rangle=\frac{N_c}{(N_c^2-4)(N_c^2-1)}\,C^{abc}d^{abc}\,\sigma(p_1,p_2,p_3)\notag\\&\quad+\frac{i}{2(N_c^2-1)}\Big[\Phi_{12}\,\delta^{(2)}(p_1-p_2)\,n(p_1,p_3)+\Phi_{23}\,\delta^{(2)}(p_2-p_3)\,n(p_1,p_2)\notag\\&\qquad\qquad-\Phi_{13}\,\delta^{(2)}(p_1-p_3)\,n(p_1,p_2)\Big]\,,\notag\\
&\langle C^{abc}\rho^a(p_1)\{\rho^b(p_2),\rho^c(p_3)\}\rangle=\frac{2N_c}{(N_c^2-4)(N_c^2-1)}\,C^{abc}d^{abc}\,\sigma(p_1,p_2,p_3)\notag\\&\quad+\frac{i}{N_c^2-1}\Big[\Phi_{12}\,\delta^{(2)}(p_1-p_2)\,n(p_1,p_3)-\Phi_{13}\,\delta^{(2)}(p_1-p_3)\,n(p_1,p_2)\Big]\,,\notag\\
&\langle C^{abc}\{\rho^a(p_1),\rho^b(p_2)\}\rho^c(p_3)\rangle=\frac{2N_c}{(N_c^2-4)(N_c^2-1)}\,C^{abc}d^{abc}\,\sigma(p_1,p_2,p_3)\notag\\&\quad+\frac{i}{N_c^2-1}\Big[\Phi_{23}\,\delta^{(2)}(p_2-p_3)\,n(p_1,p_2)-\Phi_{13}\,\delta^{(2)}(p_1-p_3)\,n(p_1,p_2)\Big]\,.
\end{align}
Every cubic row therefore has two parts after the average, and both are part of the result: the symmetrised three-$\rho$ part, proportional to $\sigma$, with the factor $t\,i$ in the equations below, $t=1$ for an ordered product and $t=2$ for a product with an anticommutator, and the contact part, two-point terms $\propto\delta^{(2)}(p_i-p_j)\,n$ with the coefficients $(\pm\frac{i}{2})$, resp.\ $(\pm i)$, of Eq.~\eqref{eq:cubicavg}. Since charge conjugation of the projectile reverses $\sigma$ and leaves $n$ unchanged, the two parts are the parts of the cubic terms that are odd and even under charge conjugation of the projectile. Section~\ref{sec:check_cubic} shows the split on a single quark and on a dipole, computed exactly in charge space.

\paragraph{Parity in $\v k$.} Write a term as $A(\v k)\langle O\rangle$, with $A$ the $c$-number factor and $O$ the colour operator. All kernels depend on $\v k$ through the phases, including the recoil phases under the $\xi$ integrals, and through the distribution~\eqref{eq:inst_xi}, which turn into their complex conjugates under $\v k\to-\v k$; the genuine two-$\rho$ rows have real prefactors and the cubic rows an explicit $i$ (in Row~I.II the $i$ is in the Fourier transform of the odd function $R^i$). Hence $A(-\v k)=A(\v k)^*$ for the genuine two-$\rho$ and $A(-\v k)=-A(\v k)^*$ for the cubic terms; for Row~II.I, which has no complex conjugate, the same holds for the sum of its two terms. With $\langle O\rangle=\mathsf S+i\mathsf C$, $\mathsf S$ the symmetrised and $\mathsf C$ the contact part, both real, a cubic term plus its complex conjugate is $\mathsf S[A(\v k)-A(-\v k)]+i\mathsf C[A(\v k)+A(-\v k)]$. The symmetrised part of the cubic terms is odd under $\v k\to-\v k$, their contact part is even, and the genuine two-$\rho$ terms are even. The part of the finite terms that is odd in $\v k$ is therefore exactly the symmetrised three-$\rho$ part, which is proportional to the projectile three-point function $\sigma$ and to C-odd structures of the target (see below). This refines the statement that the cubic, ``first saturation'' terms are odd in $\v k$~\cite{McLerran:2016snu,Kovchegov:2018jun}: it holds for their symmetrised part, configuration by configuration, but not for the commutator remainders of the ordered products.

\paragraph{Colour structures.} In all rows the colour tensors are written as sums of products of \emph{differences} of Wilson lines, as in the LO structure $\big[U^{ce'}(y')-U^{ce'}(x')\big]\big[U^{ce}(y)-U^{ce}(x)\big]$. The rewriting uses only the invariance of the structure constants, $f^{gde}U^{db}(p)U^{ec}(p)=U^{gh}(p)f^{hbc}$, the orthogonality $U^{mh}(p)U^{kh}(p)=\delta^{mk}$ and $f^{acd}f^{bcd}=N_c\delta^{ab}$, and is exact for any $N_c$; for instance,
\begin{align}\label{eq:factex}
&f^{hbc}U^{gh}(w)-f^{gde}U^{db}(x)U^{ec}(y)\notag\\&\quad=-f^{gde}\Big\{\big[U^{db}(x)-U^{db}(w)\big]U^{ec}(y)+U^{db}(w)\big[U^{ec}(y)-U^{ec}(w)\big]\Big\}\,.
\end{align}
In the averaged form the contractions $C^{abc}d^{abc}$ of the symmetrised part and $C^{abc}f^{abc}$ of the contact part are written as traces in the adjoint representation. For matrices $B$, $C$ and an orthogonal matrix $O$,
\begin{align}\label{eq:traceid}
f^{ijk}d^{lmn}\,O^{il}B^{jm}C^{kn}&=i\,{\rm Tr}\big\{T^hO^\dagger C\,D^hB^\dagger O\big\}\,,\notag\\
f^{ijk}f^{lmn}\,O^{il}B^{jm}C^{kn}&={\rm Tr}\big\{T^eO^\dagger C\,T^eB^\dagger O\big\}\,,
\end{align}
where $B$ and $C$ are the differences of Wilson lines and $O$ is a Wilson line or a product of Wilson lines, e.g.\ ${\rm Tr}\big\{T^hU^\dagger(x)\big(U(z)-U(y)\big)D^h\big(U^\dagger(x')-U^\dagger(w')\big)U(x)\big\}$. The traces with $D^h$ are imaginary; $i$ times such a trace is real and is a C-odd (odderon) structure of the target: it changes sign when the fundamental Wilson lines, $U^{ab}=2\,{\rm tr}(t^aVt^bV^\dagger)$, are replaced by their complex conjugates, it starts at third order in the target field for a weak target, and it vanishes for $SU(2)$. The traces with $T^e$ and the traces of two differences are C-even.

The extent to which each colour leg can be written as a difference is stated in Appendix~\ref{app:finite}. Every term of every row therefore vanishes for $U=1$, and for any Wilson line that does not depend on the transverse position.

\paragraph{The LO structure.} In Row~II.IV, and in the DGLAP and running-coupling terms~\eqref{eq:DGLAP}, \eqref{eq:RC} and in Eq.~\eqref{eq:LORC}, the colour factor is that of the LO spectrum, and the average is the replacement
\begin{align}\label{eq:sLOrepl}
&\big[U^{ce'}(y')-U^{ce'}(x')\big]\big[U^{ce}(y)-U^{ce}(x)\big]\rho^{e'}(x')\rho^e(x)\notag\\&\to\;\frac{n(x',x)}{N_c^2-1}\,{\rm Tr}\big\{\big(U^\dagger(y')-U^\dagger(x')\big)\big(U(y)-U(x)\big)\big\}\,.
\end{align}
In the labelling of the BFKL term ($x'\to y$, $y\to w$, $y'\to w'$) the trace is $\mathcal O_{\rm LO}$ of Eq.~\eqref{eq:OLO}.

\subsection{Recoil, soft remainders and their cancellation}\label{sec:nuance_soft}

Some finite terms exist only in combination. This subsection gives the reasons; the full analysis, with all checks, is in Appendix~\ref{app:tails}.

\paragraph{Plus prescriptions.} A term of a row has the form $\int\frac{d\xi}{\xi}f(\xi)$, with $f$ a function of $\xi$ and of all transverse coordinates. Its soft end point gives $f(0)\log(k^+/\Lambda)$, which is the entry of the row in the JIMWLK list, and the remainder $\int_0\frac{d\xi}{\xi}[f(\xi)-f(0)]$ is the plus-prescription term. The remainder converges at $\xi\to0$ at fixed transverse coordinates. Whether it is also integrable over the transverse coordinates is a separate question, and this is where the recoil enters.

\paragraph{Recoil phases.} When the measured gluon is one member of a pair, or the product of the recombination of a pair, its transverse position is tied to the $p^+$-weighted centre of the pair (Section~\ref{sec:ex_III_V}, Eq.~\eqref{eq:ex3_phase}). The phase then contains $e^{-i\xi s}$ with $s=\v k\cdot(\v z-\v w)$ in Group~I and $s=\v k\cdot(\v z-\v w')$ in Group~III. In the simplest case nothing else depends on $\xi$, $f(\xi)=g(0)e^{-i\xi s}$, and the plus-prescription term is
\begin{align}\label{eq:softrem_main}
g(0)\int_0^1\frac{d\xi}{\xi}\big(e^{-i\xi s}-1\big)=g(0)\,\mathcal E(s)\,,\qquad\mathcal E(s)\simeq-\log|s|-\gamma_E-\frac{i\pi}{2}\,{\rm sgn}(s)\quad(|s|\gg1)\,.
\end{align}
This is the \emph{soft remainder}. It has a simple meaning: the soft integral does not extend down to $\xi=0$ but is cut off at $\xi\sim1/|s|$, where the phase starts to oscillate; in terms of rapidities, the soft gluon contributes only if its rapidity is below that of the measured gluon. The logarithm $\log(k^+/\Lambda)$ of the JIMWLK list overcounts by $\log|s|$, and the soft remainder corrects for it.

\paragraph{Double logarithms.} The JIMWLK coefficient $g(0)$ falls like $T/\v z^2$ at large distance $\v z$ of the soft gluon. Multiplied by $\log(k^+/\Lambda)$ this is the familiar infrared behaviour of the JIMWLK kernel, which cancels between real and virtual terms. Multiplied by $\mathcal E\simeq-\log|\v k\cdot\v z|$ it gives $\int d^2z\,\log|\v z|/\v z^2$, which diverges like the square of a logarithm. At finite $\Lambda$ the same region produces $\pi\,T_{\rm row}\log^2(k^+/\Lambda)$ in each row with a recoil. Such double logarithms would spoil the identification of the single logarithm with the JIMWLK evolution; they cancel because $\sum_{\rm rows}T_{\rm row}=0$ (numerically $|\sum T|<10^{-7}$ against single tails of order one to six).

\paragraph{The pairing of Group~I with Group~III.} The reason for $\sum T=0$ is that the soft remainders cancel in pairs. A soft gluon can be attached to the measured gluon in two ways: it is absorbed by it in the amplitude (virtual, Group~I, recoil on $\v w$), or it is emitted by it in the conjugate amplitude and integrated over (real, Group~III, recoil on $\v w'$). The two terms have the same coefficient, the same ordered colour operator and the same transverse kernel up to $\v w\leftrightarrow\v w'$, and opposite signs. Their difference, Eq.~\eqref{eq:rmv}, is integrable: the logarithms differ only by $\log\frac{|\v k\cdot(\v z-\v w')|}{|\v k\cdot(\v z-\v w)|}=O(1/|\v z|)$, and the kernels by $O(|\v w-\v w'|/|\v z|^3)$. There are four such pairs (Table~\ref{tab:pairs}). Two consequences follow. First, the members of a pair must stand under one $\v z$ integral: neither exists alone. Second, since the colour operator, with its whole dependence on $U(\v z)$, is a common factor, the cancellation needs no assumption about the target at large distance. \emph{It is the tails that cancel, not the terms}: both members of every pair are part of the result.

\paragraph{Two terms of Row~I.III.} Row~I.III, the rescattering of the measured gluon on the valence charges, contains two terms in which the mechanism just described is hidden.

(i) \emph{A term whose pointwise limit $\Lambda\to0$ is not integrable.} One term of the row has, at fixed $\v z\neq\v w$ and $\Lambda\to0$, the kernel $\log\frac{(\v z-\v w)^2}{(\v x-\v w)^2}/(\v z-\v w)^2$. Its $\v z$ integral diverges at $\v z\to\v w$ like $-2\pi\Phi(\v w)\log^2(\varepsilon/|\v x-\v w|)$ for a cut-off $\varepsilon$ on $|\v z-\v w|$. The divergence is an artefact of taking $\Lambda\to0$ before the integral. With $\Lambda$ kept, and after the Fourier integral over $\v w$, a change of variables to the positions of the soft and of the hard gluon turns the term \emph{exactly} into
\begin{align}\label{eq:recoil_identity}
-\int_{w,z}e^{-i\v k\cdot\v w}\,K(\v x,\v w,\v z)\,\Phi(\v w)\int_{\lambda}^{1-\lambda}\frac{dt}{t}\,e^{-it\,\v k\cdot(\v z-\v w)}\,,\qquad\lambda=\frac{\Lambda}{k^+}\,,
\end{align}
i.e.\ into $\log(k^+/\Lambda)$ times a JIMWLK-list term with the Wilson line at $\v w$, plus the soft remainder with $\mathcal E(\v k\cdot(\v z-\v w))$. (The identity was checked numerically with a Gaussian regulator, to $3\times10^{-13}$.) The recoil term so obtained is the term with the kernel $\mathfrak K^{(w)}_{\rm I.III}$, the Group~I member of pair~A. Thus the limit $\Lambda\to0$ and the transverse integral do not commute term by term, and that the correct coefficient of $\log(k^+/\Lambda)$ is found only in the variable in which the recoil is a phase.

(ii) \emph{A term with a tail that no single partner cancels.} Another term has the factor $\ell=-\frac12\log\frac{(\v z-\v w)^2}{(\v x-\v w)^2}$ in place of a logarithm of the cut-off: here the soft integral is cut off by the transverse geometry, not by a phase, and nothing of it is in the JIMWLK list. At large $\v z$ it has an isotropic tail $\propto-\log|\v z|/\v z^2$. It is cancelled in two parts. Its $\log|\v z|$ part is cancelled by the soft remainder of Row~III.VI (pair~D); for this purpose $\ell$ is written as $[\ell-\mathcal E]+\mathcal E$. The remaining $1/\v z^2$ part, which depends on the direction of $\v z$, is cancelled by the tail of the \emph{other} terms of the same row, but only after the integral over $\v w$: at fixed $\v w$ those terms have a tail $\propto\varphi/\v z^2$ that is odd in the angle $\varphi$ between $\v x-\v w$ and $\v z$ and does not match, while the Fourier transforms of the two tails are equal and opposite for every direction of $\v z$, Eq.~\eqref{eq:sawid}.

\paragraph{The order of the $\v w$ and $\v z$ integrals.} In Row~I.III the integral over $\v w$ has therefore to be done before the integral over $\v z$. The integral over $\v w$ with the phase $e^{-i\v k\cdot\v w}$ \emph{is} the measurement: it projects the gluon on the momentum $\v k$, and the sum over the unmeasured gluon is to be taken at fixed $\v k$. It is a Fourier integral of a tempered function and is the same for any regulator, Eq.~\eqref{eq:FTlemma}. In the opposite order the $\v z$ integral does not exist. With this order every remaining integral converges absolutely.

\subsection{Principal values and instantaneous terms}\label{sec:nuance_pv}

\paragraph{Where vanishing denominators occur.} In the dressing of the vacuum all energy denominators are sums of positive energies and cannot vanish. In the dressing of a gluon that is already present the denominators are differences, $E_g(k)-E_n$, and vanish when the intermediate state is degenerate with the initial one. This happens in one coefficient, $\mf B_3$ (the transition of a gluon into a gluon of different longitudinal momentum on the valence charges, on the surface $p^+(\v x-\v u)^2=k^+(\v x-\v z)^2$), and therefore in one row, I.III.

\paragraph{The prescription.} These denominators are defined as principal values, Eq.~\eqref{eq:PV}. Integrals over such denominators give logarithms of absolute values, and in the kernel of Row~I.III as printed the poles of the separate terms have cancelled and all logarithms have positive arguments. A prescription $1/(E-E_n\pm i\varepsilon)$ differs from the principal value by the on-shell term $\mp i\pi\delta(E-E_n)$, and would add to Row~I.III a finite term that is not included in the result. Two facts limit what this term can be. It is confined to Row~I.III; and in the dilute limit for a point-like projectile it is absent, because there the on-shell part of $\mf B_3$ is purely imaginary relative to the LO amplitude and drops out of the interference (Section~\ref{sec:check_dilute}). Beyond this limit the dependence on the prescription is an open point (Section~\ref{sec:openpoints}).

\paragraph{Denominators that need no prescription.} The kernels of the genuine two-$\rho$ part of Row~II.I have denominators that vanish on the surface $(\v x'-\v w')^2(\v x-\v z)^2=(\v x-\v w)^2(\v x'-\v z)^2$. Here no prescription is needed: the sums of the terms $a+b$ and $c+d$ are finite and continuous across the surface, by the mechanism known from the kernels $\log(X^2Y'^2/X'^2Y^2)/(X^2Y'^2-X'^2Y^2)$ of the NLO JIMWLK Hamiltonian~\cite{Balitsky:2013fea,Kovner:2013ona}.

\paragraph{Instantaneous terms as distributions.} The opposite end of the longitudinal integral raises a different question. In Rows~III.III, III.V and~III.VI one tensor structure has no dependence on $\xi$ except through the recoil phase (the third term of Eq.~\eqref{eq:ex3_three} for Row~III.V), so that $\int_0^{\vee/k^+}d\xi$ does not converge pointwise as $\vee\to\infty$. The integral is the distribution $\int_0^\infty d\xi\,e^{-i\xi s}=\pi\delta(s)-i\mathcal P\frac1s$ in $s=\v k\cdot(\v z-\v y')$, and it becomes a number once the transverse integral over $\v y'$ is done first. Three statements make this well defined (Appendix~\ref{sec:fin_inst}): the term and its subtraction (the kernels $\mathfrak K^{(y')}_{\rm III}$, last term, and $\mathfrak K^{(z)}_{\rm III}$) must be kept together, since each alone grows like $\log\vee$; the combination vanishes at $\v y'\to\v z$, so that the distribution acts on a continuous function; and the result does not come from $p^+\sim\vee$, where the eikonal approximation would fail, but from $p^+\lesssim k^+\max\{1,1/(|\v k|\ell)\}$ with $\ell$ the shortest transverse scale. Numerically the integral cut at $\vee$ approaches the distributional value like $k^+/\vee$.

\subsection{Exact cancellations and terms that must be integrated together}\label{sec:fin_together}

We now list the cancellations that have been carried out and the terms that have to be combined under one integral.

\paragraph{Exact cancellations.} Two cancellations among the finite terms are exact, i.e.\ operator identities at the level of the integrand, and have been carried out: the finite part of Row~I.IV cancels that of Row~III.I, and the terms proportional to $\frac{11}{3}$ of Rows~I.II and~I.V, which carry the ultraviolet poles of the virtual corrections, cancel each other (Section~\ref{sec:dg_poles}). All other finite terms remain and are listed below.

\paragraph{Terms that must be integrated together.} Several of the remaining terms have a tail $\propto1/\v z^2$ at large transverse distance $\v z$ of the integrated gluon, which alone would make the $\v z$ integral diverge. These tails cancel between terms of different rows, but only the tails cancel, not the terms; none of them is dropped, and the terms concerned have to be combined under one integral. They are:
\begin{itemize}
\item \emph{Pair A:} the term of Row~I.III with the difference $\big[U(w)-U(y)\big]$, i.e.\ the term with the kernel $\mathfrak K^{(w)}_{\rm I.III}$, with Row~III.V.
\item \emph{Pair B:} the terms of the cubic part of Row~I.V with $\big[\frac1\xi\big]_+$, kernel $\mathfrak K_{\rm I.V}$, with the term of the cubic part of Row~III.III with $\big[\frac{1}{\xi(1+\xi)}\big]_+$, kernel $\mathfrak K^{(y')}_{\rm III}$.
\item \emph{Pair C:} the terms of the genuine two-$\rho$ part of Row~I.V with $\big[\frac1\xi\big]_+$ and $\big[\frac{1}{1-\xi}\big]_+$, kernel $\mathfrak K^{\rm quad}_{\rm I.V}$, with the term of the genuine two-$\rho$ part of Row~III.III with $\big[\frac{1}{\xi(1+\xi)}\big]_+$.
\item \emph{Pair D:} the part $\int_0^1d\xi\big[\frac1\xi\big]_+e^{-ik\cdot(z-w)\xi}$ of the term of Row~I.III with $\big[U(z)-U(y)\big]$, i.e.\ the last term of the kernel $\mathfrak K^{(z)}_{\rm I.III}$, with Row~III.VI.
\end{itemize}
Each pair stands under one $\v z$ integral, with $\v y'$ renamed $\v w'$ in the Group~III member (and $\v x\leftrightarrow\v y$ in pair~D). The members of a pair are a virtual term (Group~I, the integrated gluon is absorbed by the measured gluon in the amplitude, recoil phase with $\v w$) and a real term (Group~III, the integrated gluon is emitted by the measured gluon in the conjugate amplitude, recoil phase with $\v w'$) with the same coefficient, the same ordered colour operator and the same transverse kernel up to $\v w\to\v w'$; the cancellation of their tails holds for every value of $U(\v z)$ and requires no condition on the target at large distance. In Row~I.III, moreover, the remaining terms of the kernel $\mathfrak K^{(z)}_{\rm I.III}$, i.e.\ its first three groups of terms and the term with $-\frac12\log\frac{(z-w)^2}{(x-w)^2}-\int_0^1d\xi\big[\frac1\xi\big]_+e^{-ik\cdot(z-w)\xi}$, are integrable only with the $\v w$ integral, the Fourier transform at the measured momentum, done before the $\v z$ integral. Finally, Row~II.IV, the genuine two-$\rho$ part of Row~II.I and the genuine two-$\rho$ part of Row~III.III have to be integrated together over the position of the pair of final-state gluons: in the region in which both gluons are far from the projectile their tails cancel with the weights $1:1:-2$ (Section~\ref{sec:dg_tail}). The origin of these statements, the closed forms of the longitudinal integrals involved and the checks are collected in Appendix~\ref{app:tails}.

\paragraph{Row I.II: the momentum-space kernel.} The cubic part of Row~I.II is given as a Fourier integral of a momentum-space kernel, $\mathfrak K_{\rm I.II}$ of Eq.~\eqref{eq:K-rm_I_II}. This kernel is the antisymmetrised one, $\frac12\big[R^i(\v q,\v p)/\v p^2-R^i(\v q,\v q-\v p)/(\v q-\v p)^2\big]$: each of its two parts alone behaves as $\log(1/r)/r^2$ at $\v p\to\v q$ ($r=|\v p-\v q|$) and at $\v p\to0$ ($r=|\v p|$), which is not integrable over $d^2p$, while in the antisymmetrised combination these terms cancel and the kernel is $O(\log(1/r)/r)$, which is integrable. The antisymmetrised form is equivalent to the form with $R^i(\v q,\v p)/\v p^2$ alone, since the colour structure of the row is antisymmetric under $\v x\leftrightarrow\v y$, which under the phases $e^{iq\cdot(w-y)}e^{ip\cdot(y-x)}$ is $\v p\to\v q-\v p$; the antisymmetrised form makes explicit the value that the scaleless rule $\int d^2p/\v p^2=0$ of dimensional regularisation assigns to these points.

\paragraph{Row I.V: the collinear point.} In the genuine two-$\rho$ part of Row~I.V the kernel $\mathfrak K^{\rm quad}_{\rm I.V}$ is singular at the collinear point $\v z\to\v x$. There the term is finite only together with its $\frac{11}{6}$ line, the second term of Eq.~\eqref{eq:S-I.V-2o}, and with a rotation-invariant treatment of the point $\v z=\v x$, i.e.\ with the angular average around this point taken first: the sum of the two behaves as a pure second harmonic in the angle of $\v z-\v x$, divided by $(\v z-\v x)^2$. The rotation-invariant treatment is the one consistent with the $\frac{11}{6}$ term, which comes from a dimensionally regularised, rotation-invariant integral.

\subsection{The finite terms in colour-singlet form}\label{sec:fin_rows}

We now list the finite terms, averaged over the projectile according to Eqs.~\eqref{eq:A1n}--\eqref{eq:cubicavg}. The corresponding expressions before the average are given in Appendix~\ref{app:finite}; the two lists have the same prefactors, kernels and phases. Blocks of transverse and longitudinal kernels that occur in several equations are given names, $\mathfrak K$ with the row as a subscript. A kernel depends on the transverse integration variables of the equation in which it stands, on $\v k$ and, where it stands under a $\xi$ integral, on $\xi$; a plus prescription inside a kernel acts on everything in the term that depends on $\xi$, including a recoil phase written in front of the kernel, according to Eq.~\eqref{eq:plus_xi}. In the contact terms the $\delta$-functions are to be used to set the charge positions equal in the kernel as well.

\begingroup
% the generated equations carry \small; the list is set in \footnotesize
\let\small\footnotesize

\subsubsection{Group I}

\paragraph{Row I.II.} With the kernel
\begingroup\small
\begin{align}\label{eq:K-rm_I_II}
&\mathfrak{K}_{\rm I.II}\equiv \int d^{2}p\, d^{2}q\;e^{iq\cdot(w-y)}\, e^{ip\cdot(y-x)}\, \frac{(x'-w')^{i}}{(x'-w')^{2}\,\mathbf{q}^{2}}\, \frac{1}{2}\bigg[\frac{R^{i}(\mathbf{q},\mathbf{p})}{\mathbf{p}^{2}} -\frac{R^{i}(\mathbf{q},\mathbf{q}-\mathbf{p})}{(\mathbf{q}-\mathbf{p})^{2}}\bigg]\,,
\end{align}
\endgroup

where
\begin{align}\label{eq:Rdef}
&R^{i}(\mathbf{q},\mathbf{p})=-2q^{i}\log\left(\frac{\mathbf{p}^{2}}{\mathbf{q}^{2}}\right) +\frac{(2\mathbf{p}^{2}-2\,\mathbf{q}\cdot\mathbf{p}+\mathbf{q}^{2})\,q^{i}+\mathbf{q}^{2}p^{i}}{(\mathbf{q}-\mathbf{p})^{2}}\, \log\left[\frac{\mathbf{q}^{2}}{(\mathbf{q}-\mathbf{p})^{2}}\right]\,,
\end{align}

the symmetrised three-$\rho$ part is
\begingroup\small
\begin{align}\label{eq:S-I.II-3}
&\Big\langle\frac{d^{3}N^{\rm I.II}_{{\rm sym}}}{d^{3}k}\Big\rangle=-\frac{1}{(2\pi)^{3}}\, \frac{g^{4}}{64\pi^{6}}\, \frac{1}{k^{+}} \int_{w,w'}e^{-ik\cdot(w-w')} \int_{x,y,x'} \mathfrak{K}_{\rm I.II} \times \frac{N_{c}}{(N_{c}^{2}-4)(N_{c}^{2}-1)}\, \Big(+2i\Big)\notag\\
&\quad \times \sigma(x',x,y)\Big[ -{\rm Tr}\Big\{T^{h}U^{\dagger}(y)\big(U(x)-U(w)\big)D^{h}\big(U^{\dagger}(x')-U^{\dagger}(w')\big)U(y)\Big\}\notag\\
&\quad\quad -{\rm Tr}\Big\{T^{h}U^{\dagger}(w)\big(U(x')-U(w')\big)D^{h}\big(U^{\dagger}(y)-U^{\dagger}(w)\big)U(w)\Big\}\Big] +{\rm c.c.}
\end{align}
\endgroup

the contact part is
\begingroup\small
\begin{align}\label{eq:S-I.II-2}
&\Big\langle\frac{d^{3}N^{\rm I.II}_{{\rm con}}}{d^{3}k}\Big\rangle=-\frac{1}{(2\pi)^{3}}\, \frac{g^{4}}{64\pi^{6}}\, \frac{1}{k^{+}} \int_{w,w'}e^{-ik\cdot(w-w')} \int_{x,y,x'} \mathfrak{K}_{\rm I.II} \times \frac{1}{N_{c}^{2}-1}\bigg[\Big(+i\Big)\delta^{(2)}(x'-x)\notag\\
&\quad\quad \times n(x,y)\Big[ -{\rm Tr}\Big\{T^{e}U^{\dagger}(y)\big(U(x)-U(w)\big)T^{e}\big(U^{\dagger}(x)-U^{\dagger}(w')\big)U(y)\Big\}\notag\\
&\quad\quad\quad -{\rm Tr}\Big\{T^{e}U^{\dagger}(w)\big(U(y)-U(w)\big)T^{e}\big(U^{\dagger}(x)-U^{\dagger}(w')\big)U(w)\Big\}\Big]\notag\\
&\quad\quad +\Big(-i\Big)\delta^{(2)}(x'-y)\, n(y,x)\Big[ -{\rm Tr}\Big\{T^{e}U^{\dagger}(y)\big(U(x)-U(w)\big)T^{e}\big(U^{\dagger}(y)-U^{\dagger}(w')\big)U(y)\Big\}\notag\\
&\quad\quad\quad -{\rm Tr}\Big\{T^{e}U^{\dagger}(w)\big(U(y)-U(w)\big)T^{e}\big(U^{\dagger}(y)-U^{\dagger}(w')\big)U(w)\Big\}\Big]\bigg] +{\rm c.c.}
\end{align}
\endgroup

and the genuine two-$\rho$ term, which is proportional to the LO structure, is
\begingroup\small
\begin{align}\label{eq:S-I.II-2o}
&\Big\langle\frac{d^{3}N^{\rm I.II}_{{\rm quad}}}{d^{3}k}\Big\rangle=+\frac{1}{(2\pi)^{3}}\, \frac{g^{4}\,N_{c}}{32\pi^{4}}\, \frac{1}{k^{+}}\Big(-\frac{67}{9}+\frac{\pi^{2}}{3}\Big) \int_{w,w',x,x'} e^{-ik\cdot(w-w')}\;\Kv(x'-w')\cdot\Kv(x-w)\notag\\
&\quad \times\frac{n(x',x)}{N_{c}^{2}-1}\Big[ {\rm Tr}\Big\{\big(U^{\dagger}(x')-U^{\dagger}(w')\big)\big(U(w)-U(x)\big)\Big\}\Big] +{\rm c.c.}
\end{align}
\endgroup

The constant $-\frac{67}{9}+\frac{\pi^2}{3}$ is the one left in this row by the cancellation of the $\frac{11}{3}$ terms against Row~I.V; it is a property of this grouping of the terms and not a constant of the result (Section~\ref{sec:lit_constant}).

\paragraph{Row I.III.} The kernels are
\begingroup\small
\begin{align}\label{eq:K-z_rm_I_III}
&\mathfrak{K}^{(z)}_{\rm I.III}\equiv \bigg[\frac{\Kv(x'-w')\cdot\Kv(y-z)}{2\,(z-w)^{2}} +\Kv(x'-w')\cdot\Kv(z-w)\, \Kv(y-z)\cdot\Kv(x-z)\notag\\
&\quad\quad -\Kv(x'-w')\cdot\Kv(x-w)\, \Kv(y-z)\cdot\Kv(z-w) -\Kv(x'-w')\cdot\Kv(x-w)\notag\\
&\quad\quad \times \Kv(y-z)\cdot\Kv(x-z)\bigg]\log\frac{(x-z)^{2}}{(x-w)^{2}} +\frac{\arccos\left(\frac{(x-w)^{2}-(x-w)\cdot (x-z)}{\sqrt{(x-w)^{2}\,(z-w)^{2}}}\right)}{\sqrt{(x-w)^{2}\,(x-z)^{2}-\left[(x-w)\cdot (x-z)\right]^{2}}\,(z-w)^{2}}\notag\\
&\quad \times \bigg[\Kv(x'-w')\cdot(z-w)\, \Kv(y-z)\cdot(x-w) -\frac{(x-w)\cdot(x-z)}{(x-z)^{2}}\, \Kv(x'-w')\cdot(z-w)\notag\\
&\quad\quad \times \Kv(y-z)\cdot(x-z) +\Kv(x'-w')\cdot(x-z)\, \Kv(y-z)\cdot(z-w)\notag\\
&\quad\quad -\frac{(x-w)\cdot(x-z)}{(x-w)^{2}}\, \Kv(x'-w')\cdot(x-w)\, \Kv(y-z)\cdot(z-w)\bigg]\notag\\
&\quad -\frac{1}{2\,(z-w)^{2}}\log\frac{(z-w)^{2}}{(x-w)^{2}}\, \Kv(x'-w')\cdot(z-w)\, \Kv(y-z)\cdot\Kv(x-z)\notag\\
&\quad +\bigg[-\frac{1}{2}\log\frac{(z-w)^{2}}{(x-w)^{2}} - \int_{0}^{1}d\xi\Big[\frac{1}{\xi}\Big]_{+}e^{-ik\cdot(z-w)\,\xi}\bigg]\Kv(y-z)\cdot\Kv(w-z)\, \Kv(x'-w')\cdot\Kv(x-w)\notag\\
&\quad + \int_{0}^{1}d\xi\Big[\frac{1}{\xi}\Big]_{+}e^{-ik\cdot(z-w)\,\xi}\, \Kv(y-z)\cdot\Kv(w-z)\, \Kv(x'-w')\cdot\Kv(x-w)\,,
\end{align}
\endgroup

and
\begingroup\small
\begin{align}\label{eq:K-w_rm_I_III}
&\mathfrak{K}^{(w)}_{\rm I.III}\equiv - \int_{0}^{1}d\xi\Big[\frac{1}{\xi}\Big]_{+}e^{-ik\cdot(z-w)\,\xi}\, \Kv(x-z)\cdot\Kv(w-z)\, \Kv(x'-w')\cdot\Kv(y-w)\,.
\end{align}
\endgroup

The last two terms of $\mathfrak K^{(z)}_{\rm I.III}$ are written separately, although the integrals over $\xi$ cancel between them, because they are integrated differently: the last term is the Group~I member of pair~D, while the term before it belongs with the rest of the kernel (Section~\ref{sec:fin_together}). The symmetrised three-$\rho$ part of the row is
\begingroup\small
\begin{align}\label{eq:S-I.III-3}
&\Big\langle\frac{d^{3}N^{\rm I.III}_{{\rm sym}}}{d^{3}k}\Big\rangle=\frac{1}{(2\pi)^{3}}\, \frac{ig^{4}}{8\pi^{5}}\, \frac{1}{k^{+}} \int_{w,w'}e^{-ik\cdot(w-w')} \int_{x,y,z,x'}\Bigg[ \mathfrak{K}^{(z)}_{\rm I.III} \times \frac{N_{c}}{(N_{c}^{2}-4)(N_{c}^{2}-1)}\, \Big(+i\Big)\notag\\
&\quad\quad \times \sigma(x',x,y)\Big[ {\rm Tr}\Big\{T^{h}U^{\dagger}(x)\big(U(z)-U(y)\big)D^{h}\big(U^{\dagger}(x')-U^{\dagger}(w')\big)U(x)\Big\}\Big]\notag\\
&\quad\quad + \mathfrak{K}^{(w)}_{\rm I.III} \times \frac{N_{c}}{(N_{c}^{2}-4)(N_{c}^{2}-1)}\, \Big(+i\Big)\, \sigma(x',x,y)\Big[\notag\\
&\quad\quad\quad {\rm Tr}\Big\{T^{h}U^{\dagger}(x)\big(U(w)-U(y)\big)D^{h}\big(U^{\dagger}(x')-U^{\dagger}(w')\big)U(x)\Big\}\Big]\Bigg] +{\rm c.c.}
\end{align}
\endgroup

and the contact part is
\begingroup\small
\begin{align}\label{eq:S-I.III-2}
&\Big\langle\frac{d^{3}N^{\rm I.III}_{{\rm con}}}{d^{3}k}\Big\rangle=\frac{1}{(2\pi)^{3}}\, \frac{ig^{4}}{8\pi^{5}}\, \frac{1}{k^{+}} \int_{w,w'}e^{-ik\cdot(w-w')} \int_{x,y,z,x'}\Bigg[ \mathfrak{K}^{(z)}_{\rm I.III} \times \frac{1}{N_{c}^{2}-1}\bigg[\Big(+\frac{1}{2}i\Big)\delta^{(2)}(x'-x)\notag\\
&\quad\quad\quad \times n(x,y)\Big[ -{\rm Tr}\Big\{T^{e}U^{\dagger}(x)\big(U(z)-U(y)\big)T^{e}\big(U^{\dagger}(x)-U^{\dagger}(w')\big)U(x)\Big\}\Big]\notag\\
&\quad\quad\quad +\Big(+\frac{1}{2}i\Big)\delta^{(2)}(y-x)\, n(x,x')\Big[ -{\rm Tr}\Big\{T^{e}U^{\dagger}(x)\big(U(z)-U(x)\big)T^{e}\big(U^{\dagger}(x')-U^{\dagger}(w')\big)U(x)\Big\}\Big]\notag\\
&\quad\quad\quad +\Big(-\frac{1}{2}i\Big)\delta^{(2)}(x'-y)\, n(y,x)\Big[ -{\rm Tr}\Big\{T^{e}U^{\dagger}(x)\big(U(z)-U(y)\big)T^{e}\big(U^{\dagger}(y)-U^{\dagger}(w')\big)U(x)\Big\}\Big]\bigg]\notag\\
&\quad\quad + \mathfrak{K}^{(w)}_{\rm I.III} \times \frac{1}{N_{c}^{2}-1}\bigg[\Big(+\frac{1}{2}i\Big)\delta^{(2)}(x'-x)\, n(x,y)\Big[\notag\\
&\quad\quad\quad\quad -{\rm Tr}\Big\{T^{e}U^{\dagger}(x)\big(U(w)-U(y)\big)T^{e}\big(U^{\dagger}(x)-U^{\dagger}(w')\big)U(x)\Big\}\Big]\notag\\
&\quad\quad\quad +\Big(+\frac{1}{2}i\Big)\delta^{(2)}(y-x)\, n(x,x')\Big[ -{\rm Tr}\Big\{T^{e}U^{\dagger}(x)\big(U(w)-U(x)\big)T^{e}\big(U^{\dagger}(x')-U^{\dagger}(w')\big)U(x)\Big\}\Big]\notag\\
&\quad\quad\quad +\Big(-\frac{1}{2}i\Big)\delta^{(2)}(x'-y)\, n(y,x)\Big[\notag\\
&\quad\quad\quad\quad -{\rm Tr}\Big\{T^{e}U^{\dagger}(x)\big(U(w)-U(y)\big)T^{e}\big(U^{\dagger}(y)-U^{\dagger}(w')\big)U(x)\Big\}\Big]\bigg]\Bigg] +{\rm c.c.}
\end{align}
\endgroup

The row has no genuine two-$\rho$ term.

\paragraph{Row I.V.} The kernel of the cubic part is
\begingroup\small
\begin{align}\label{eq:K-rm_I_V}
&\mathfrak{K}_{\rm I.V}\equiv e^{-ik\cdot(x-z)\,\xi}\, \Kv(x'-w')\cdot\Kv(z-x)\, \Kv(y-z)\cdot\Kv(y'-x)\notag\\
&\quad -\Big[\frac{1}{\xi}\Big]_{+}e^{-ik\cdot(x-z)\,\xi}\, \Kv(x'-w')\cdot\Kv(y-z)\, \Kv(z-x)\cdot\Kv(y'-x)\notag\\
&\quad -\Big[\frac{1}{\xi}\Big]_{+}e^{-ik\cdot(x-z)}e^{-ik\cdot(z-x)\,\xi}\, \Kv(x'-w')\cdot\Kv(y'-x)\, \Kv(z-x)\cdot\Kv(y-z)\,.
\end{align}
\endgroup

The symmetrised three-$\rho$ part is
\begingroup\small
\begin{align}\label{eq:S-I.V-3}
&\Big\langle\frac{d^{3}N^{\rm I.V}_{{\rm sym}}}{d^{3}k}\Big\rangle=\frac{1}{(2\pi)^{3}}\, \frac{ig^{4}}{16\pi^{4}}\, \frac{1}{k^{+}} \int_{0}^{1}\frac{d\xi}{2\pi} \int_{x,z,y,y',x',w'}e^{-ik\cdot(z-w')}\notag\\
&\quad \times \mathfrak{K}_{\rm I.V} \times \frac{N_{c}}{(N_{c}^{2}-4)(N_{c}^{2}-1)}\, \Big(+2i\Big)\, \sigma(x',y,y')\Big[\notag\\
&\quad\quad {\rm Tr}\Big\{T^{h}U^{\dagger}(z)\big(U(x)-U(y')\big)D^{h}\big(U^{\dagger}(x')-U^{\dagger}(w')\big)U(z)\Big\}\notag\\
&\quad\quad +{\rm Tr}\Big\{T^{h}U^{\dagger}(y')\big(U(x')-U(w')\big)D^{h}\big(U^{\dagger}(z)-U^{\dagger}(y)\big)U(y')\Big\}\Big] +{\rm c.c.}
\end{align}
\endgroup

and the contact part is
\begingroup\small
\begin{align}\label{eq:S-I.V-2}
&\Big\langle\frac{d^{3}N^{\rm I.V}_{{\rm con}}}{d^{3}k}\Big\rangle=\frac{1}{(2\pi)^{3}}\, \frac{ig^{4}}{16\pi^{4}}\, \frac{1}{k^{+}} \int_{0}^{1}\frac{d\xi}{2\pi} \int_{x,z,y,y',x',w'}e^{-ik\cdot(z-w')}\notag\\
&\quad \times \mathfrak{K}_{\rm I.V} \times \frac{1}{N_{c}^{2}-1}\bigg[\Big(+i\Big)\delta^{(2)}(x'-y)\, n(y,y')\Big[\notag\\
&\quad\quad\quad -{\rm Tr}\Big\{T^{e}U^{\dagger}(z)\big(U(x)-U(y')\big)T^{e}\big(U^{\dagger}(y)-U^{\dagger}(w')\big)U(z)\Big\}\notag\\
&\quad\quad\quad -{\rm Tr}\Big\{T^{e}U^{\dagger}(y')\big(U(z)-U(y)\big)T^{e}\big(U^{\dagger}(y)-U^{\dagger}(w')\big)U(y')\Big\}\Big]\notag\\
&\quad\quad +\Big(-i\Big)\delta^{(2)}(x'-y')\, n(y',y)\Big[ -{\rm Tr}\Big\{T^{e}U^{\dagger}(z)\big(U(x)-U(y')\big)T^{e}\big(U^{\dagger}(y')-U^{\dagger}(w')\big)U(z)\Big\}\notag\\
&\quad\quad\quad -{\rm Tr}\Big\{T^{e}U^{\dagger}(y')\big(U(z)-U(y)\big)T^{e}\big(U^{\dagger}(y')-U^{\dagger}(w')\big)U(y')\Big\}\Big]\bigg] +{\rm c.c.}
\end{align}
\endgroup

The genuine two-$\rho$ part has the kernel
\begingroup\small
\begin{align}\label{eq:K-rm_quad_rm_I_V}
&\mathfrak{K}^{\rm quad}_{\rm I.V}\equiv \frac{e^{-ik\cdot(x-z)\,\xi}}{\xi\,(y-x)^{2}+(1-\xi)\,(y-z)^{2}}\bigg[\xi(1-\xi)\, \frac{(y-x)^{2}-(y-z)^{2}}{(x-z)^{2}}\, \Kv(x'-w')\cdot\Kv(z-x)\notag\\
&\quad\quad +(1-\xi)\, \Kv(x'-w')\cdot\Kv(z-x)\Big[(y-z)\cdot\Kv(x-z) +\frac{1}{2}(y-z)\cdot\Kv(y-x)\notag\\
&\quad\quad\quad -\frac{(y-x)^{2}-(y-z)^{2}}{2\,(x-z)^{2}}\Big] +\xi\, \Kv(x'-w')\cdot\Kv(z-x)\Big[(y-x)\cdot\Kv(x-z) -\frac{1}{2}(y-x)\cdot\Kv(y-z)\notag\\
&\quad\quad\quad -\frac{(y-x)^{2}-(y-z)^{2}}{2\,(x-z)^{2}}\Big] +\frac{(y-x)^{2}}{(y-z)^{2}}\Big[-\frac{(y-z)\cdot\Kv(x'-w')}{(x-z)^{2}}\notag\\
&\quad\quad\quad +\frac{1}{2}\, (y-z)\cdot\Kv(x'-w')\, \Kv(z-x)\cdot\Kv(y-x)\Big]\notag\\
&\quad\quad +\frac{(y-z)^{2}}{(y-x)^{2}}\Big[-\frac{(y-x)\cdot\Kv(x'-w')}{(x-z)^{2}} -\frac{1}{2}\, (y-x)\cdot\Kv(x'-w')\, \Kv(z-x)\cdot\Kv(y-z)\Big]\notag\\
&\quad\quad -(y-x)\cdot\Kv(z-x)\, \Kv(x-z)\cdot\Kv(x'-w') +\frac{1}{2}\, (y-x)\cdot\Kv(z-x)\, \Kv(y-z)\cdot\Kv(x'-w')\notag\\
&\quad\quad -(y-z)\cdot\Kv(z-x)\, \Kv(x-z)\cdot\Kv(x'-w') -\frac{1}{2}\, \Kv(y-x)\cdot\Kv(x'-w')\, (y-z)\cdot\Kv(z-x)\bigg]\notag\\
&\quad -\Big[\frac{1}{\xi}\Big]_{+}\frac{e^{-ik\cdot(x-z)\,\xi}}{(y-z)^{2}}\Big[-\frac{(y-z)\cdot\Kv(x'-w')}{(x-z)^{2}} +\frac{1}{2}\, (y-z)\cdot\Kv(x'-w')\, \Kv(z-x)\cdot\Kv(y-x)\Big]\notag\\
&\quad -\Big[\frac{1}{1-\xi}\Big]_{+}\frac{e^{-ik\cdot(x-z)\,\xi}}{(y-x)^{2}}\Big[-\frac{(y-x)\cdot\Kv(x'-w')}{(x-z)^{2}} -\frac{1}{2}\, (y-x)\cdot\Kv(x'-w')\, \Kv(z-x)\cdot\Kv(y-z)\Big]\,,
\end{align}
\endgroup

and reads
\begingroup\small
\begin{align}\label{eq:S-I.V-2o}
&\Big\langle\frac{d^{3}N^{\rm I.V}_{{\rm quad}}}{d^{3}k}\Big\rangle=\frac{1}{(2\pi)^{3}}\, \frac{g^{4}}{8\pi^{4}}\, \frac{1}{k^{+}} \int_{0}^{1}\frac{d\xi}{2\pi} \int_{x,y,z,x',w'}e^{-ik\cdot(z-w')}\notag\\
&\quad \times \mathfrak{K}^{\rm quad}_{\rm I.V} \times\frac{n(x',y)}{N_{c}^{2}-1}\Big[ -{\rm Tr}\Big\{T^{e}U^{\dagger}(x)\big(U(x')-U(w')\big)T^{e}\big(U^{\dagger}(z)-U^{\dagger}(x)\big)U(x)\Big\}\notag\\
&\quad\quad +N_{c}{\rm Tr}\Big\{\big(U^{\dagger}(x')-U^{\dagger}(w')\big)\big(U(y)-U(x)\big)\Big\}\Big] +\frac{1}{(2\pi)^{3}}\, \frac{g^{4}N_{c}}{16\pi^{5}}\, \frac{1}{k^{+}} \int_{x,y,z,x',w'}e^{-ik\cdot(z-w')}\notag\\
&\quad \times e^{-ik\cdot(x-z)}\, \frac{11}{6}\, \Kv(x'-w')\cdot\Kv(y-x)\, \frac{(y-x)\cdot\Kv(y-z)}{(x-z)^{2}}\notag\\
&\quad \times\frac{n(x',y)}{N_{c}^{2}-1}\Big[ {\rm Tr}\Big\{\big(U^{\dagger}(x')-U^{\dagger}(w')\big)\big(U(y)-U(x)\big)\Big\}\Big]\, +{\rm c.c.}\, ,
\end{align}
\endgroup

The second term, with $\frac{11}{6}$, is the one required at the collinear point $\v z\to\v x$ (Section~\ref{sec:fin_together}).

\subsubsection{Group II}

\paragraph{Row II.I.} The kernels of the cubic part are
\begingroup\small
\begin{align}\label{eq:K-1_rm_II_I}
&\mathfrak{K}^{(1)}_{\rm II.I}\equiv -\bigg[\frac{\Kv(x'-w')\cdot\Kv(y-z)}{2\,(z-w)^{2}} +\Kv(x'-w')\cdot\Kv(z-w)\, \Kv(y-z)\cdot\Kv(x-z)\notag\\
&\quad\quad -\Kv(x'-w')\cdot\Kv(x-w)\, \Kv(y-z)\cdot\Kv(x-z) -\Kv(x'-w')\cdot\Kv(x-w)\notag\\
&\quad\quad \times \Kv(y-z)\cdot\Kv(z-w)\bigg]\log\frac{(x-z)^{2}}{(x-w)^{2}}\,,
\end{align}
\endgroup

and
\begingroup\small
\begin{align}\label{eq:K-2_rm_II_I}
&\mathfrak{K}^{(2)}_{\rm II.I}\equiv \bigg[\frac{\Kv(x'-w)\cdot\Kv(y-z)}{2\,(z-w')^{2}} +\Kv(x'-w)\cdot\Kv(z-w')\, \Kv(y-z)\cdot\Kv(x-z)\notag\\
&\quad\quad -\Kv(x'-w)\cdot\Kv(x-w')\, \Kv(y-z)\cdot\Kv(x-z) -\Kv(x'-w)\cdot\Kv(x-w')\notag\\
&\quad\quad \times \Kv(y-z)\cdot\Kv(z-w')\bigg]\log\frac{(x-z)^{2}}{(x-w')^{2}}\,,
\end{align}
\endgroup

which is $-\mathfrak K^{(1)}_{\rm II.I}$ with $\v w\leftrightarrow\v w'$. The symmetrised three-$\rho$ part is
\begingroup\small
\begin{align}\label{eq:S-II.I-3}
&\Big\langle\frac{d^{3}N^{\rm II.I}_{{\rm sym}}}{d^{3}k}\Big\rangle=\frac{1}{(2\pi)^{3}}\, \frac{ig^{4}}{16\pi^{5}}\, \frac{1}{k^{+}} \int_{x,y,z,w,x',w'}e^{-ik\cdot(w-w')}\Bigg[ \mathfrak{K}^{(1)}_{\rm II.I} \times \frac{N_{c}}{(N_{c}^{2}-4)(N_{c}^{2}-1)}\, \Big(+2i\Big)\notag\\
&\quad\quad \times \sigma(x',y,x)\Big[ -{\rm Tr}\Big\{T^{h}U^{\dagger}(z)\big(U(w')-U(x')\big)D^{h}\big(U^{\dagger}(x)-U^{\dagger}(z)\big)U(z)\Big\}\notag\\
&\quad\quad\quad +{\rm Tr}\Big\{T^{h}U^{\dagger}(z)\big(U(z)-U(y)\big)D^{h}U^{\dagger}(x')\big(U(w)-U(z)\big)\Big\}\notag\\
&\quad\quad\quad -{\rm Tr}\Big\{T^{h}U^{\dagger}(x')\big(U(x)-U(z)\big)D^{h}\big(U^{\dagger}(z)-U^{\dagger}(y)\big)U(x')\Big\}\Big]\notag\\
&\quad\quad + \mathfrak{K}^{(2)}_{\rm II.I} \times \frac{N_{c}}{(N_{c}^{2}-4)(N_{c}^{2}-1)}\, \Big(+2i\Big)\, \sigma(x,x',y)\Big[\notag\\
&\quad\quad\quad -{\rm Tr}\Big\{T^{h}U^{\dagger}(z)\big(U(w)-U(x')\big)D^{h}\big(U^{\dagger}(x)-U^{\dagger}(z)\big)U(z)\Big\}\notag\\
&\quad\quad\quad +{\rm Tr}\Big\{T^{h}U^{\dagger}(z)\big(U(z)-U(y)\big)D^{h}U^{\dagger}(x')\big(U(w')-U(z)\big)\Big\}\notag\\
&\quad\quad\quad -{\rm Tr}\Big\{T^{h}U^{\dagger}(x')\big(U(x)-U(z)\big)D^{h}\big(U^{\dagger}(z)-U^{\dagger}(y)\big)U(x')\Big\}\Big]\Bigg]
\end{align}
\endgroup

and the contact part is
\begingroup\small
\begin{align}\label{eq:S-II.I-2}
&\Big\langle\frac{d^{3}N^{\rm II.I}_{{\rm con}}}{d^{3}k}\Big\rangle=\frac{1}{(2\pi)^{3}}\, \frac{ig^{4}}{16\pi^{5}}\, \frac{1}{k^{+}} \int_{x,y,z,w,x',w'}e^{-ik\cdot(w-w')}\Bigg[ \mathfrak{K}^{(1)}_{\rm II.I} \times \frac{1}{N_{c}^{2}-1}\bigg[\Big(+i\Big)\delta^{(2)}(y-x)\notag\\
&\quad\quad\quad \times n(x,x')\Big[ N_{c}{\rm Tr}\Big\{\big(U^{\dagger}(w')-U^{\dagger}(x')\big)\big(U(w)-U(z)\big)\Big\}\notag\\
&\quad\quad\quad\quad -{\rm Tr}\Big\{T^{e}U^{\dagger}(z)\big(U(x)-U(z)\big)T^{e}\big(U^{\dagger}(w')-U^{\dagger}(x')\big)U(z)\Big\}\notag\\
&\quad\quad\quad\quad +{\rm Tr}\Big\{T^{e}U^{\dagger}(z)\big(U(z)-U(x)\big)T^{e}U^{\dagger}(x')\big(U(w)-U(z)\big)\Big\}\notag\\
&\quad\quad\quad\quad -{\rm Tr}\Big\{T^{e}U^{\dagger}(x')\big(U(x)-U(z)\big)T^{e}\big(U^{\dagger}(z)-U^{\dagger}(x)\big)U(x')\Big\}\Big]\notag\\
&\quad\quad\quad +\Big(-i\Big)\delta^{(2)}(x'-x)\, n(x,y)\Big[ N_{c}{\rm Tr}\Big\{\big(U^{\dagger}(w')-U^{\dagger}(x)\big)\big(U(w)-U(z)\big)\Big\}\notag\\
&\quad\quad\quad\quad -{\rm Tr}\Big\{T^{e}U^{\dagger}(z)\big(U(x)-U(z)\big)T^{e}\big(U^{\dagger}(w')-U^{\dagger}(x)\big)U(z)\Big\}\notag\\
&\quad\quad\quad\quad +{\rm Tr}\Big\{T^{e}U^{\dagger}(z)\big(U(z)-U(y)\big)T^{e}U^{\dagger}(x)\big(U(w)-U(z)\big)\Big\}\notag\\
&\quad\quad\quad\quad -{\rm Tr}\Big\{T^{e}U^{\dagger}(x)\big(U(x)-U(z)\big)T^{e}\big(U^{\dagger}(z)-U^{\dagger}(y)\big)U(x)\Big\}\Big]\bigg]\notag\\
&\quad\quad + \mathfrak{K}^{(2)}_{\rm II.I} \times \frac{1}{N_{c}^{2}-1}\bigg[\Big(+i\Big)\delta^{(2)}(x'-x)\, n(x,y)\Big[ N_{c}{\rm Tr}\Big\{\big(U^{\dagger}(w)-U^{\dagger}(x)\big)\big(U(w')-U(z)\big)\Big\}\notag\\
&\quad\quad\quad\quad -{\rm Tr}\Big\{T^{e}U^{\dagger}(z)\big(U(x)-U(z)\big)T^{e}\big(U^{\dagger}(w)-U^{\dagger}(x)\big)U(z)\Big\}\notag\\
&\quad\quad\quad\quad +{\rm Tr}\Big\{T^{e}U^{\dagger}(z)\big(U(z)-U(y)\big)T^{e}U^{\dagger}(x)\big(U(w')-U(z)\big)\Big\}\notag\\
&\quad\quad\quad\quad -{\rm Tr}\Big\{T^{e}U^{\dagger}(x)\big(U(x)-U(z)\big)T^{e}\big(U^{\dagger}(z)-U^{\dagger}(y)\big)U(x)\Big\}\Big]\notag\\
&\quad\quad\quad +\Big(-i\Big)\delta^{(2)}(y-x)\, n(x,x')\Big[ N_{c}{\rm Tr}\Big\{\big(U^{\dagger}(w)-U^{\dagger}(x')\big)\big(U(w')-U(z)\big)\Big\}\notag\\
&\quad\quad\quad\quad -{\rm Tr}\Big\{T^{e}U^{\dagger}(z)\big(U(x)-U(z)\big)T^{e}\big(U^{\dagger}(w)-U^{\dagger}(x')\big)U(z)\Big\}\notag\\
&\quad\quad\quad\quad +{\rm Tr}\Big\{T^{e}U^{\dagger}(z)\big(U(z)-U(x)\big)T^{e}U^{\dagger}(x')\big(U(w')-U(z)\big)\Big\}\notag\\
&\quad\quad\quad\quad -{\rm Tr}\Big\{T^{e}U^{\dagger}(x')\big(U(x)-U(z)\big)T^{e}\big(U^{\dagger}(z)-U^{\dagger}(x)\big)U(x')\Big\}\Big]\bigg]\Bigg]
\end{align}
\endgroup

Group~II has no complex conjugate. The genuine two-$\rho$ part consists of five terms with the same colour factor. With the kernels
\begingroup\small
\begin{align}\label{eq:K-rm_quad_a_rm_II_I}
&\mathfrak{K}^{{\rm quad},a}_{\rm II.I}\equiv \frac{1}{\big[(x'-w')^{2}\,(x-z)^{2}-(x-w)^{2}\,(x'-z)^{2}\big]}\, \log\left[\frac{(x-z)^{2}}{(x'-z)^{2}}\right]\notag\\
&\quad \times\Bigg[(x-w)\cdot\Kv(z-w')\, (x'-z)\cdot\Kv(z-w) -\frac{1}{2}\, (x-w)\cdot\Kv(x'-w')\, \Kv(z-w)\cdot(x'-z)\notag\\
&\quad\quad +\frac{1}{2}\, \Kv(x-z)\cdot(x'-z)\, (x-w)\cdot\Kv(z-w') -\frac{1}{4}\, (x'-z)\cdot\Kv(x-z)\, \Kv(x'-w')\cdot(x-w)\notag\\
&\quad\quad +(x'-w')\cdot\Kv(z-w)\, (x-z)\cdot\Kv(z-w') +\frac{1}{2}\, (x'-w')\cdot\Kv(z-w)\, (x-z)\cdot\Kv(x'-z)\notag\\
&\quad\quad -\frac{1}{2}\, (x-z)\cdot\Kv(z-w')\, (x'-w')\cdot\Kv(x-w) -\frac{1}{4}\, (x'-w')\cdot\Kv(x-w)\, (x-z)\cdot\Kv(x'-z)\Bigg]\,,
\end{align}
\endgroup

\begingroup\small
\begin{align}\label{eq:K-rm_quad_b_rm_II_I}
&\mathfrak{K}^{{\rm quad},b}_{\rm II.I}\equiv \frac{1}{\big[(x'-w')^{2}\,(x-z)^{2}-(x-w)^{2}\,(x'-z)^{2}\big]}\, \log\left[\frac{(x-w)^{2}}{(x'-w')^{2}}\right]\notag\\
&\quad \times\Bigg[(x-w)\cdot\Kv(z-w')\, (x'-z)\cdot\Kv(z-w) -\frac{1}{2}\, (x-w)\cdot\Kv(x'-w')\, \Kv(z-w)\cdot(x'-z)\notag\\
&\quad\quad +\frac{1}{2}\, \Kv(x-z)\cdot(x'-z)\, (x-w)\cdot\Kv(z-w') -\frac{1}{4}\, (x'-z)\cdot\Kv(x-z)\, \Kv(x'-w')\cdot(x-w)\notag\\
&\quad\quad +(x'-w')\cdot\Kv(z-w)\, (x-z)\cdot\Kv(z-w') +\frac{1}{2}\, (x'-w')\cdot\Kv(z-w)\, (x-z)\cdot\Kv(x'-z)\notag\\
&\quad\quad -\frac{1}{2}\, (x-z)\cdot\Kv(z-w')\, (x'-w')\cdot\Kv(x-w) -\frac{1}{4}\, (x'-w')\cdot\Kv(x-w)\, (x-z)\cdot\Kv(x'-z)\Bigg]\,,
\end{align}
\endgroup

\begingroup\small
\begin{align}\label{eq:K-rm_quad_c_rm_II_I}
&\mathfrak{K}^{{\rm quad},c}_{\rm II.I}\equiv \frac{1}{\big[(x'-w')^{2}\,(x-z)^{2}-(x-w)^{2}\,(x'-z)^{2}\big]}\, \log\left[\frac{(x-z)^{2}}{(x-w)^{2}}\right]\notag\\
&\quad \times\Bigg[-\frac{(x-z)^{2}}{(x-w)^{2}}\bigg[\frac{(x'-w')\cdot(x-w)\,(z-w)\cdot(z-w')}{(z-w)^{2}\,(z-w')^{2}} +\frac{(x'-w')\cdot(x-w)\,(x'-z)\cdot(z-w)}{2\,(z-w)^{2}\,(x'-z)^{2}}\notag\\
&\quad\quad\quad +\frac{(x-w)\cdot(x'-w')\,(x-z)\cdot(z-w')}{2\,(x-z)^{2}\,(z-w')^{2}} +\frac{(x-z)\cdot(x'-z)\,(x'-w')\cdot(x-w)}{4\,(x-z)^{2}\,(x'-z)^{2}}\bigg]\notag\\
&\quad\quad -\frac{(x-w)^{2}}{(x-z)^{2}}\bigg[\frac{(x'-z)\cdot(x-z)\,(z-w)\cdot(z-w')}{(z-w)^{2}\,(z-w')^{2}} -\frac{(x'-z)\cdot(x-z)\,(x'-w')\cdot(z-w)}{2\,(z-w)^{2}\,(x'-w')^{2}}\notag\\
&\quad\quad\quad -\frac{(x-w)\cdot(z-w')\,(x-z)\cdot(x'-z)}{2\,(x-w)^{2}\,(z-w')^{2}} +\frac{(x-z)\cdot(x'-z)\,(x'-w')\cdot(x-w)}{4\,(x-w)^{2}\,(x'-w')^{2}}\bigg]\notag\\
&\quad\quad +\frac{1}{(x-z)^{2}-(x-w)^{2}}\bigg\{(x-z)^{2}\bigg[\frac{(x-w)\cdot(z-w)\,\big[(x'-z)^{2}-(x'-w')^{2}\big]}{2\,(z-w)^{2}\,(z-w')^{2}}\notag\\
&\quad\quad\quad\quad +\frac{(x-z)\cdot(x-w)\,\big[(x'-z)^{2}-(x'-w')^{2}\big]}{4\,(x-z)^{2}\,(z-w')^{2}} +\frac{(x'-w')\cdot(z-w')\,\big[(x-z)^{2}-(x-w)^{2}\big]}{2\,(z-w)^{2}\,(z-w')^{2}}\notag\\
&\quad\quad\quad\quad +\frac{(x'-z)\cdot(x'-w')\,\big[(x-z)^{2}-(x-w)^{2}\big]}{4\,(z-w)^{2}\,(x'-z)^{2}}\bigg]\notag\\
&\quad\quad\quad -(x-w)^{2}\bigg[\frac{(x-z)\cdot(z-w)\,\big[(x'-z)^{2}-(x'-w')^{2}\big]}{2\,(z-w)^{2}\,(z-w')^{2}}\notag\\
&\quad\quad\quad\quad -\frac{(x-z)\cdot(x-w)\,\big[(x'-z)^{2}-(x'-w')^{2}\big]}{4\,(x-w)^{2}\,(z-w')^{2}} +\frac{(x'-z)\cdot(z-w')\,\big[(x-z)^{2}-(x-w)^{2}\big]}{2\,(z-w)^{2}\,(z-w')^{2}}\notag\\
&\quad\quad\quad\quad -\frac{(x'-z)\cdot(x'-w')\,\big[(x-z)^{2}-(x-w)^{2}\big]}{4\,(z-w)^{2}\,(x'-w')^{2}}\bigg]\bigg\}\notag\\
&\quad\quad -\frac{1}{2}\, \frac{(x-z)^{2}\,(x-w)^{2}\,\big[(x'-z)^{2}-(x'-w')^{2}\big]}{(z-w)^{2}\,(z-w')^{2}\,\big[(x-z)^{2}-(x-w)^{2}\big]}\Bigg]\,,
\end{align}
\endgroup

\begingroup\small
\begin{align}\label{eq:K-rm_quad_d_rm_II_I}
&\mathfrak{K}^{{\rm quad},d}_{\rm II.I}\equiv \frac{1}{\big[(x'-w')^{2}\,(x-z)^{2}-(x-w)^{2}\,(x'-z)^{2}\big]}\, \log\left[\frac{(x'-z)^{2}}{(x'-w')^{2}}\right]\notag\\
&\quad \times\Bigg[\frac{(x'-z)^{2}}{(x'-w')^{2}}\bigg[(x'-w')\cdot(x-w)\, \Kv(z-w)\cdot\Kv(z-w')\notag\\
&\quad\quad\quad +\frac{1}{2}\, (x'-w')\cdot(x-w)\, \Kv(x'-z)\cdot\Kv(z-w) +\frac{1}{2}\, (x-w)\cdot(x'-w')\, \Kv(x-z)\cdot\Kv(z-w')\notag\\
&\quad\quad\quad +\frac{1}{4}\, \Kv(x-z)\cdot\Kv(x'-z)\, (x'-w')\cdot(x-w)\bigg]\notag\\
&\quad\quad +\frac{(x'-w')^{2}}{(x'-z)^{2}}\bigg[(x'-z)\cdot(x-z)\, \Kv(z-w)\cdot\Kv(z-w')\notag\\
&\quad\quad\quad -\frac{1}{2}\, (x'-z)\cdot(x-z)\, \Kv(x'-w')\cdot\Kv(z-w) -\frac{1}{2}\, \Kv(x-w)\cdot\Kv(z-w')\, (x-z)\cdot(x'-z)\notag\\
&\quad\quad\quad +\frac{1}{4}\, (x-z)\cdot(x'-z)\, \Kv(x'-w')\cdot\Kv(x-w)\bigg]\notag\\
&\quad\quad +\frac{1}{(x'-w')^{2}-(x'-z)^{2}}\bigg\{(x'-z)^{2}\bigg[\frac{(x-w)\cdot\Kv(z-w)\,\big[(x'-z)^{2}-(x'-w')^{2}\big]}{2\,(z-w')^{2}}\notag\\
&\quad\quad\quad\quad +\frac{(x-w)\cdot\Kv(x-z)\,\big[(x'-z)^{2}-(x'-w')^{2}\big]}{4\,(z-w')^{2}}\notag\\
&\quad\quad\quad\quad +\frac{(x'-w')\cdot\Kv(z-w')\,\big[(x-z)^{2}-(x-w)^{2}\big]}{2\,(z-w)^{2}}\notag\\
&\quad\quad\quad\quad +\frac{(x'-w')\cdot\Kv(x'-z)\,\big[(x-z)^{2}-(x-w)^{2}\big]}{4\,(z-w)^{2}}\bigg]\notag\\
&\quad\quad\quad -(x'-w')^{2}\bigg[\frac{(x-z)\cdot\Kv(z-w)\,\big[(x'-z)^{2}-(x'-w')^{2}\big]}{2\,(z-w')^{2}}\notag\\
&\quad\quad\quad\quad -\frac{(x-z)\cdot\Kv(x-w)\,\big[(x'-z)^{2}-(x'-w')^{2}\big]}{4\,(z-w')^{2}}\notag\\
&\quad\quad\quad\quad +\frac{(x'-z)\cdot\Kv(z-w')\,\big[(x-z)^{2}-(x-w)^{2}\big]}{2\,(z-w)^{2}}\notag\\
&\quad\quad\quad\quad -\frac{(x'-z)\cdot\Kv(x'-w')\,\big[(x-z)^{2}-(x-w)^{2}\big]}{4\,(z-w)^{2}}\bigg]\bigg\}\notag\\
&\quad\quad +\frac{1}{2}\, \frac{(x'-w')^{2}\,(x'-z)^{2}\,\big[(x-z)^{2}-(x-w)^{2}\big]}{(z-w)^{2}\,(z-w')^{2}\,\big[(x'-z)^{2}-(x'-w')^{2}\big]}\Bigg]\,,
\end{align}
\endgroup

the five terms are
\begingroup\small
\begin{align}\label{eq:S-II.I-2a}
&\Big\langle\frac{d^{3}N^{\rm II.I}_{{\rm quad},a}}{d^{3}k}\Big\rangle=\frac{1}{(2\pi)^{3}}\, \frac{g^{4}}{8\pi^{5}}\, \frac{1}{k^{+}} \int_{z,w,w',x,x'}e^{-ik\cdot(w-w')}\, \mathfrak{K}^{{\rm quad},a}_{\rm II.I} \times\frac{n(x',x)}{N_{c}^{2}-1}\Big[\notag\\
&\quad\quad N_{c}{\rm Tr}\Big\{\big(U^{\dagger}(w')-U^{\dagger}(z)\big)\big(U(w)-U(z)\big)\Big\} +N_{c}{\rm Tr}\Big\{\big(U^{\dagger}(x')-U^{\dagger}(z)\big)\big(U(x)-U(z)\big)\Big\}\notag\\
&\quad\quad -{\rm Tr}\Big\{T^{e}U^{\dagger}(z)\big(U(x)-U(z)\big)T^{e}\big(U^{\dagger}(w')-U^{\dagger}(z)\big)U(z)\Big\}\notag\\
&\quad\quad -{\rm Tr}\Big\{T^{e}U^{\dagger}(z)\big(U(x')-U(z)\big)T^{e}\big(U^{\dagger}(w)-U^{\dagger}(z)\big)U(z)\Big\}\Big]
\end{align}
\endgroup

\begingroup\small
\begin{align}\label{eq:S-II.I-2b}
&\Big\langle\frac{d^{3}N^{\rm II.I}_{{\rm quad},b}}{d^{3}k}\Big\rangle=-\frac{1}{(2\pi)^{3}}\, \frac{g^{4}}{8\pi^{5}}\, \frac{1}{k^{+}} \int_{z,w,w',x,x'}e^{-ik\cdot(w-w')}\, \mathfrak{K}^{{\rm quad},b}_{\rm II.I} \times\frac{n(x',x)}{N_{c}^{2}-1}\Big[\notag\\
&\quad\quad N_{c}{\rm Tr}\Big\{\big(U^{\dagger}(w')-U^{\dagger}(z)\big)\big(U(w)-U(z)\big)\Big\} +N_{c}{\rm Tr}\Big\{\big(U^{\dagger}(x')-U^{\dagger}(z)\big)\big(U(x)-U(z)\big)\Big\}\notag\\
&\quad\quad -{\rm Tr}\Big\{T^{e}U^{\dagger}(z)\big(U(x)-U(z)\big)T^{e}\big(U^{\dagger}(w')-U^{\dagger}(z)\big)U(z)\Big\}\notag\\
&\quad\quad -{\rm Tr}\Big\{T^{e}U^{\dagger}(z)\big(U(x')-U(z)\big)T^{e}\big(U^{\dagger}(w)-U^{\dagger}(z)\big)U(z)\Big\}\Big]
\end{align}
\endgroup

\begingroup\small
\begin{align}\label{eq:S-II.I-2c}
&\Big\langle\frac{d^{3}N^{\rm II.I}_{{\rm quad},c}}{d^{3}k}\Big\rangle=\frac{1}{(2\pi)^{3}}\, \frac{g^{4}}{8\pi^{5}}\, \frac{1}{k^{+}} \int_{z,w,w',x,x'}e^{-ik\cdot(w-w')}\, \mathfrak{K}^{{\rm quad},c}_{\rm II.I} \times\frac{n(x',x)}{N_{c}^{2}-1}\Big[\notag\\
&\quad\quad N_{c}{\rm Tr}\Big\{\big(U^{\dagger}(w')-U^{\dagger}(z)\big)\big(U(w)-U(z)\big)\Big\} +N_{c}{\rm Tr}\Big\{\big(U^{\dagger}(x')-U^{\dagger}(z)\big)\big(U(x)-U(z)\big)\Big\}\notag\\
&\quad\quad -{\rm Tr}\Big\{T^{e}U^{\dagger}(z)\big(U(x)-U(z)\big)T^{e}\big(U^{\dagger}(w')-U^{\dagger}(z)\big)U(z)\Big\}\notag\\
&\quad\quad -{\rm Tr}\Big\{T^{e}U^{\dagger}(z)\big(U(x')-U(z)\big)T^{e}\big(U^{\dagger}(w)-U^{\dagger}(z)\big)U(z)\Big\}\Big]
\end{align}
\endgroup

\begingroup\small
\begin{align}\label{eq:S-II.I-2d}
&\Big\langle\frac{d^{3}N^{\rm II.I}_{{\rm quad},d}}{d^{3}k}\Big\rangle=\frac{1}{(2\pi)^{3}}\, \frac{g^{4}}{8\pi^{5}}\, \frac{1}{k^{+}} \int_{z,w,w',x,x'}e^{-ik\cdot(w-w')}\, \mathfrak{K}^{{\rm quad},d}_{\rm II.I} \times\frac{n(x',x)}{N_{c}^{2}-1}\Big[\notag\\
&\quad\quad N_{c}{\rm Tr}\Big\{\big(U^{\dagger}(w')-U^{\dagger}(z)\big)\big(U(w)-U(z)\big)\Big\} +N_{c}{\rm Tr}\Big\{\big(U^{\dagger}(x')-U^{\dagger}(z)\big)\big(U(x)-U(z)\big)\Big\}\notag\\
&\quad\quad -{\rm Tr}\Big\{T^{e}U^{\dagger}(z)\big(U(x)-U(z)\big)T^{e}\big(U^{\dagger}(w')-U^{\dagger}(z)\big)U(z)\Big\}\notag\\
&\quad\quad -{\rm Tr}\Big\{T^{e}U^{\dagger}(z)\big(U(x')-U(z)\big)T^{e}\big(U^{\dagger}(w)-U^{\dagger}(z)\big)U(z)\Big\}\Big]
\end{align}
\endgroup

\begingroup\small
\begin{align}\label{eq:S-II.I-2e}
&\Big\langle\frac{d^{3}N^{\rm II.I}_{{\rm quad},e}}{d^{3}k}\Big\rangle=-\frac{1}{(2\pi)^{3}}\, \frac{g^{4}}{16\pi^{5}}\, \frac{1}{k^{+}} \int_{z,w,w',x,x'}e^{-ik\cdot(w-w')}\, \frac{1}{(z-w)^{2}\,(z-w')^{2}}\notag\\
&\quad \times\frac{n(x',x)}{N_{c}^{2}-1}\Big[ N_{c}{\rm Tr}\Big\{\big(U^{\dagger}(w')-U^{\dagger}(z)\big)\big(U(w)-U(z)\big)\Big\}\notag\\
&\quad\quad +N_{c}{\rm Tr}\Big\{\big(U^{\dagger}(x')-U^{\dagger}(z)\big)\big(U(x)-U(z)\big)\Big\}\notag\\
&\quad\quad -{\rm Tr}\Big\{T^{e}U^{\dagger}(z)\big(U(x)-U(z)\big)T^{e}\big(U^{\dagger}(w')-U^{\dagger}(z)\big)U(z)\Big\}\notag\\
&\quad\quad -{\rm Tr}\Big\{T^{e}U^{\dagger}(z)\big(U(x')-U(z)\big)T^{e}\big(U^{\dagger}(w)-U^{\dagger}(z)\big)U(z)\Big\}\Big]
\end{align}
\endgroup

The poles of the kernels $\mathfrak K^{{\rm quad},a}_{\rm II.I},\dots,\mathfrak K^{{\rm quad},d}_{\rm II.I}$ on the surface $(x'-w')^2(x-z)^2=(x-w)^2(x'-z)^2$ cancel in the sums of the terms $a$ and $b$ and of the terms $c$ and $d$.

\paragraph{Row II.IV.} With $\xi$ of Eq.~\eqref{eq:xi_defIV},
\begingroup\small
\begin{align}\label{eq:S-II.IVrem}
&\Big\langle\frac{d^{3}N^{\rm II.IV}_{{\rm fin}}}{d^{3}k}\Big\rangle=+\frac{1}{(2\pi)^{3}}\, \frac{g^{4}\,N_{c}}{4\pi^{3}}\, \frac{1}{k^{+}} \int_{k^{+}/\vee}^{1}\frac{d\xi}{2\pi}\, \frac{1-\xi}{\xi} \int_{x,y,x',y'}e^{-ik\cdot(y-y')\frac{1}{\xi}}\Big[\Kv(x-y)\cdot\Kv(x'-y')\notag\\
&\quad\quad -2\, \Kv(x-y)\cdot(y'-y)\, \Kv(x'-y')\cdot\Kv(y'-y)\Big] \times\frac{n(x',x)}{N_{c}^{2}-1}\notag\\
&\quad \times {\rm Tr}\Big\{\big(U^{\dagger}(y')-U^{\dagger}(x')\big)\big(U(y)-U(x)\big)\Big\} -\frac{1}{(2\pi)^{3}}\, \frac{g^{4}\,N_{c}}{2\pi^{3}}\, \frac{1}{k^{2}}\, \frac{1}{k^{+}}\notag\\
&\quad \times \int_{x,x',z}\Kv(x-z)\cdot\Kv(x'-z)\frac{n(x',x)}{N_{c}^{2}-1}\, {\rm Tr}\Big\{\big(U^{\dagger}(z)-U^{\dagger}(x')\big)\big(U(z)-U(x)\big)\Big\}
\end{align}
\endgroup

The first term is the finite term of the real splitting that is left beside the DGLAP term~\eqref{eq:DGLAP}; its tensor structure averages to zero over the direction of $\v y'-\v y$. After the integral over $\v y-\v y'$ its integrand is $O(\xi)$ at $\xi\to0$, so that the limit $\vee\to\infty$ exists; at fixed coordinates the $\xi$ integral converges conditionally and gives the same result. The second term is the remnant~\eqref{eq:remnant} of the subtraction at $\xi\to0$ in $\widetilde{\mathcal P}_{gg}$.

\subsubsection{Group III}

\paragraph{Row III.II.} The kernel is
\begingroup\small
\begin{align}\label{eq:K-rm_III_II}
&\mathfrak{K}_{\rm III.II}\equiv \bigg[\frac{\Kv(x'-w')\cdot\Kv(y-z)}{2\,(z-w)^{2}} +\Kv(x'-w')\cdot\Kv(z-w)\, \Kv(y-z)\cdot\Kv(x-z)\notag\\
&\quad\quad -\Kv(x'-w')\cdot\Kv(x-w)\, \Kv(y-z)\cdot\Kv(x-z) -\Kv(x'-w')\cdot\Kv(x-w)\notag\\
&\quad\quad \times \Kv(y-z)\cdot\Kv(z-w)\bigg]\log\frac{(x-z)^{2}}{(x-w)^{2}}\,,
\end{align}
\endgroup

so that $\mathfrak K^{(1)}_{\rm II.I}=-\mathfrak K_{\rm III.II}$. The symmetrised three-$\rho$ part is
\begingroup\small
\begin{align}\label{eq:S-III.II-3}
&\Big\langle\frac{d^{3}N^{\rm III.II}_{{\rm sym}}}{d^{3}k}\Big\rangle=-\frac{1}{(2\pi)^{3}}\, \frac{ig^{4}}{8\pi^{5}}\, \frac{1}{k^{+}} \int_{x,y,x',z,w,w'}e^{-ik\cdot(w-w')}\notag\\
&\quad \times \mathfrak{K}_{\rm III.II} \times \frac{N_{c}}{(N_{c}^{2}-4)(N_{c}^{2}-1)}\, \Big(+i\Big)\, \sigma(y,x',x)\Big[\notag\\
&\quad\quad -{\rm Tr}\Big\{T^{h}U^{\dagger}(z)\big(U(z)-U(y)\big)D^{h}U^{\dagger}(x')\big(U(w)-U(z)\big)\Big\}\notag\\
&\quad\quad +{\rm Tr}\Big\{T^{h}U^{\dagger}(x')\big(U(x)-U(z)\big)D^{h}\big(U^{\dagger}(z)-U^{\dagger}(y)\big)U(x')\Big\}\Big] +{\rm c.c.}
\end{align}
\endgroup

and the contact part is
\begingroup\small
\begin{align}\label{eq:S-III.II-2}
&\Big\langle\frac{d^{3}N^{\rm III.II}_{{\rm con}}}{d^{3}k}\Big\rangle=-\frac{1}{(2\pi)^{3}}\, \frac{ig^{4}}{8\pi^{5}}\, \frac{1}{k^{+}} \int_{x,y,x',z,w,w'}e^{-ik\cdot(w-w')}\notag\\
&\quad \times \mathfrak{K}_{\rm III.II} \times \frac{1}{N_{c}^{2}-1}\bigg[\Big(+\frac{1}{2}i\Big)\delta^{(2)}(x'-y)\, n(y,x)\Big[\notag\\
&\quad\quad\quad {\rm Tr}\Big\{T^{e}U^{\dagger}(z)\big(U(z)-U(y)\big)T^{e}U^{\dagger}(y)\big(U(w)-U(z)\big)\Big\}\notag\\
&\quad\quad\quad -{\rm Tr}\Big\{T^{e}U^{\dagger}(y)\big(U(x)-U(z)\big)T^{e}\big(U^{\dagger}(z)-U^{\dagger}(y)\big)U(y)\Big\}\Big]\notag\\
&\quad\quad +\Big(+\frac{1}{2}i\Big)\delta^{(2)}(x'-x)\, n(x,y)\Big[ {\rm Tr}\Big\{T^{e}U^{\dagger}(z)\big(U(z)-U(y)\big)T^{e}U^{\dagger}(x)\big(U(w)-U(z)\big)\Big\}\notag\\
&\quad\quad\quad -{\rm Tr}\Big\{T^{e}U^{\dagger}(x)\big(U(x)-U(z)\big)T^{e}\big(U^{\dagger}(z)-U^{\dagger}(y)\big)U(x)\Big\}\Big]\notag\\
&\quad\quad +\Big(-\frac{1}{2}i\Big)\delta^{(2)}(y-x)\, n(x,x')\Big[ {\rm Tr}\Big\{T^{e}U^{\dagger}(z)\big(U(z)-U(x)\big)T^{e}U^{\dagger}(x')\big(U(w)-U(z)\big)\Big\}\notag\\
&\quad\quad\quad -{\rm Tr}\Big\{T^{e}U^{\dagger}(x')\big(U(x)-U(z)\big)T^{e}\big(U^{\dagger}(z)-U^{\dagger}(x)\big)U(x')\Big\}\Big]\bigg] +{\rm c.c.}
\end{align}
\endgroup


\paragraph{Row III.III.} Rows~III.III, III.V and~III.VI have the same kernels in their cubic parts,
\begingroup\small
\begin{align}\label{eq:K-y_rm_III}
&\mathfrak{K}^{(y')}_{\rm III}\equiv \frac{1}{1+\xi}\Big[\Kv(y'-z)\cdot\Kv(x'-y')\, \Kv(y-w)\cdot\Kv(x-z)\notag\\
&\quad\quad -\Kv(y'-z)\cdot\Kv(x-z)\, \Kv(x'-y')\cdot\Kv(y-w)\Big] -\Big[\frac{1}{\xi\,(1+\xi)}\Big]_{+}\Kv(y'-z)\cdot\Kv(x-z)\notag\\
&\quad \times \Kv(x'-y')\cdot\Kv(y-w) -\Kv(y'-z)\cdot\Kv(y-w)\, \Kv(x-z)\cdot\Kv(x'-y')\,,
\end{align}
\endgroup

and
\begingroup\small
\begin{align}\label{eq:K-z_rm_III}
&\mathfrak{K}^{(z)}_{\rm III}\equiv \Kv(y'-z)\cdot\Kv(y-w)\, \Kv(x-z)\cdot\Kv(x'-z)\,.
\end{align}
\endgroup

The last term of $\mathfrak K^{(y')}_{\rm III}$ and the kernel $\mathfrak K^{(z)}_{\rm III}$ are the instantaneous terms, defined by Eq.~\eqref{eq:inst_xi}. The symmetrised three-$\rho$ part of Row~III.III is
\begingroup\small
\begin{align}\label{eq:S-III.III-3}
&\Big\langle\frac{d^{3}N^{\rm III.III}_{{\rm sym}}}{d^{3}k}\Big\rangle=\frac{1}{(2\pi)^{3}}\, \frac{ig^{4}}{8\pi^{4}}\, \frac{1}{k^{+}} \int_{0}^{\infty}\frac{d\xi}{2\pi} \int_{w,x',y',z,y,x}e^{-ik\cdot(w-y')}e^{-ik\cdot(z-y')\,\xi}\Bigg[\notag\\
&\quad\quad \mathfrak{K}^{(y')}_{\rm III} \times \frac{N_{c}}{(N_{c}^{2}-4)(N_{c}^{2}-1)}\, \Big(+2i\Big)\, \sigma(x',y,x)\Big[\notag\\
&\quad\quad\quad {\rm Tr}\Big\{T^{h}U^{\dagger}(w)\big(U(z)-U(x)\big)D^{h}\big(U^{\dagger}(y')-U^{\dagger}(x')\big)U(w)\Big\}\notag\\
&\quad\quad\quad +{\rm Tr}\Big\{T^{h}U^{\dagger}(x)\big(U(y')-U(x')\big)D^{h}\big(U^{\dagger}(w)-U^{\dagger}(y)\big)U(x)\Big\}\Big]\notag\\
&\quad\quad + \mathfrak{K}^{(z)}_{\rm III} \times \frac{N_{c}}{(N_{c}^{2}-4)(N_{c}^{2}-1)}\, \Big(+2i\Big)\, \sigma(x',y,x)\Big[\notag\\
&\quad\quad\quad {\rm Tr}\Big\{T^{h}U^{\dagger}(w)\big(U(z)-U(x)\big)D^{h}\big(U^{\dagger}(z)-U^{\dagger}(x')\big)U(w)\Big\}\notag\\
&\quad\quad\quad +{\rm Tr}\Big\{T^{h}U^{\dagger}(x)\big(U(z)-U(x')\big)D^{h}\big(U^{\dagger}(w)-U^{\dagger}(y)\big)U(x)\Big\}\Big]\Bigg] +{\rm c.c.}
\end{align}
\endgroup

and the contact part is
\begingroup\small
\begin{align}\label{eq:S-III.III-2}
&\Big\langle\frac{d^{3}N^{\rm III.III}_{{\rm con}}}{d^{3}k}\Big\rangle=\frac{1}{(2\pi)^{3}}\, \frac{ig^{4}}{8\pi^{4}}\, \frac{1}{k^{+}} \int_{0}^{\infty}\frac{d\xi}{2\pi} \int_{w,x',y',z,y,x}e^{-ik\cdot(w-y')}e^{-ik\cdot(z-y')\,\xi}\Bigg[\notag\\
&\quad\quad \mathfrak{K}^{(y')}_{\rm III} \times \frac{1}{N_{c}^{2}-1}\bigg[\Big(+i\Big)\delta^{(2)}(x'-y)\, n(y,x)\Big[\notag\\
&\quad\quad\quad\quad -{\rm Tr}\Big\{T^{e}U^{\dagger}(w)\big(U(z)-U(x)\big)T^{e}\big(U^{\dagger}(y')-U^{\dagger}(y)\big)U(w)\Big\}\notag\\
&\quad\quad\quad\quad -{\rm Tr}\Big\{T^{e}U^{\dagger}(x)\big(U(w)-U(y)\big)T^{e}\big(U^{\dagger}(y')-U^{\dagger}(y)\big)U(x)\Big\}\Big]\notag\\
&\quad\quad\quad +\Big(-i\Big)\delta^{(2)}(x'-x)\, n(x,y)\Big[ -{\rm Tr}\Big\{T^{e}U^{\dagger}(w)\big(U(z)-U(x)\big)T^{e}\big(U^{\dagger}(y')-U^{\dagger}(x)\big)U(w)\Big\}\notag\\
&\quad\quad\quad\quad -{\rm Tr}\Big\{T^{e}U^{\dagger}(x)\big(U(w)-U(y)\big)T^{e}\big(U^{\dagger}(y')-U^{\dagger}(x)\big)U(x)\Big\}\Big]\bigg]\notag\\
&\quad\quad + \mathfrak{K}^{(z)}_{\rm III} \times \frac{1}{N_{c}^{2}-1}\bigg[\Big(+i\Big)\delta^{(2)}(x'-y)\, n(y,x)\Big[\notag\\
&\quad\quad\quad\quad -{\rm Tr}\Big\{T^{e}U^{\dagger}(w)\big(U(z)-U(x)\big)T^{e}\big(U^{\dagger}(z)-U^{\dagger}(y)\big)U(w)\Big\}\notag\\
&\quad\quad\quad\quad -{\rm Tr}\Big\{T^{e}U^{\dagger}(x)\big(U(w)-U(y)\big)T^{e}\big(U^{\dagger}(z)-U^{\dagger}(y)\big)U(x)\Big\}\Big]\notag\\
&\quad\quad\quad +\Big(-i\Big)\delta^{(2)}(x'-x)\, n(x,y)\Big[ -{\rm Tr}\Big\{T^{e}U^{\dagger}(w)\big(U(z)-U(x)\big)T^{e}\big(U^{\dagger}(z)-U^{\dagger}(x)\big)U(w)\Big\}\notag\\
&\quad\quad\quad\quad -{\rm Tr}\Big\{T^{e}U^{\dagger}(x)\big(U(w)-U(y)\big)T^{e}\big(U^{\dagger}(z)-U^{\dagger}(x)\big)U(x)\Big\}\Big]\bigg]\Bigg] +{\rm c.c.}
\end{align}
\endgroup

The genuine two-$\rho$ part is
\begingroup\small
\begin{align}\label{eq:S-III.III-2o}
&\Big\langle\frac{d^{3}N^{\rm III.III}_{{\rm quad}}}{d^{3}k}\Big\rangle=\frac{1}{(2\pi)^{3}}\, \frac{g^{4}}{4\pi^{4}}\, \frac{1}{k^{+}} \int_{0}^{\infty}\frac{d\xi}{2\pi} \int_{w,x',y',z,x}e^{-ik\cdot(w-y')}e^{-ik\cdot(z-y')\,\xi}\, \frac{1}{\xi\,(x-z)^{2}+(x-w)^{2}}\notag\\
&\quad \times\Bigg\{\Kv(y'-z)\cdot\Kv(x'-y')\bigg[\frac{\xi\big[(x-z)^{2}-(x-w)^{2}\big]}{(z-w)^{2}\,(1+\xi)^{2}}\notag\\
&\quad\quad\quad +\frac{\xi}{1+\xi}\Big((x-z)\cdot\Kv(z-w)-\frac{1}{2}(x-z)\cdot\Kv(x-w)-\frac{(x-z)^{2}-(x-w)^{2}}{2\,(z-w)^{2}}\Big)\notag\\
&\quad\quad\quad +\frac{1}{1+\xi}\Big((x-w)\cdot\Kv(z-w)+\frac{1}{2}(x-w)\cdot\Kv(x-z)-\frac{(x-z)^{2}-(x-w)^{2}}{2\,(z-w)^{2}}\Big)\bigg]\notag\\
&\quad\quad +\frac{(x-z)^{2}-(x-w)^{2}}{(1+\xi)\,(x-w)^{2}}\Big[(x-w)\cdot\Kv(x'-y')\, \Kv(z-w)\cdot\Kv(y'-z)\notag\\
&\quad\quad\quad +\frac{1}{2}\, \Kv(x-z)\cdot\Kv(y'-z)\, (x-w)\cdot\Kv(x'-y')\Big]\notag\\
&\quad\quad -\Big[(x-z)\cdot\Kv(y'-z)\, \Kv(z-w)\cdot\Kv(x'-y') -\frac{1}{2}(x-z)\cdot\Kv(y'-z)\, \Kv(x-w)\cdot\Kv(x'-y')\notag\\
&\quad\quad\quad +(x-w)\cdot\Kv(y'-z)\, \Kv(z-w)\cdot\Kv(x'-y') +\frac{1}{2}\, \Kv(x-z)\cdot\Kv(x'-y')\notag\\
&\quad\quad\quad \times (x-w)\cdot\Kv(y'-z)\Big]\Bigg\} \times\frac{n(x',x)}{N_{c}^{2}-1}\Big[ N_{c}{\rm Tr}\Big\{\big(U^{\dagger}(y')-U^{\dagger}(x')\big)\big(U(x)-U(z)\big)\Big\}\notag\\
&\quad\quad -{\rm Tr}\Big\{T^{e}U^{\dagger}(z)\big(U(w)-U(z)\big)T^{e}\big(U^{\dagger}(y')-U^{\dagger}(x')\big)U(z)\Big\}\Big]\notag\\
&\quad -\frac{1}{(2\pi)^{3}}\, \frac{g^{4}}{4\pi^{4}}\, \frac{1}{k^{+}} \int_{0}^{\infty}\frac{d\xi}{2\pi}\Big[\frac{1}{\xi\,(1+\xi)}\Big]_{+} \int_{w,x',y',z,x}e^{-ik\cdot(w-y')}e^{-ik\cdot(z-y')\,\xi}\, \frac{1}{(x-w)^{2}}\notag\\
&\quad \times\Big[(x-w)\cdot\Kv(x'-y')\, \Kv(z-w)\cdot\Kv(y'-z) +\frac{1}{2}\, \Kv(x-z)\cdot\Kv(y'-z)\, (x-w)\cdot\Kv(x'-y')\Big]\notag\\
&\quad \times\frac{n(x',x)}{N_{c}^{2}-1}\Big[ N_{c}{\rm Tr}\Big\{\big(U^{\dagger}(y')-U^{\dagger}(x')\big)\big(U(x)-U(z)\big)\Big\}\notag\\
&\quad\quad -{\rm Tr}\Big\{T^{e}U^{\dagger}(z)\big(U(w)-U(z)\big)T^{e}\big(U^{\dagger}(y')-U^{\dagger}(x')\big)U(z)\Big\}\Big]\notag\\
&\quad -\frac{1}{(2\pi)^{3}}\, \frac{g^{4}}{4\pi^{4}}\, \frac{1}{k^{+}} \int_{0}^{\infty}\frac{d\xi}{2\pi} \int_{w,x',y',z,x}e^{-ik\cdot(w-y')}e^{-ik\cdot(z-y')\,\xi}\, \frac{\xi}{\xi\,(x-z)^{2}+(x-w)^{2}}\notag\\
&\quad \times\bigg\{\Big[(x-z)\cdot\Kv(x'-y')\, \Kv(z-w)\cdot\Kv(y'-z) -\frac{1}{2}(x-z)\cdot\Kv(x'-y')\, \Kv(x-w)\cdot\Kv(y'-z)\Big]\notag\\
&\quad\quad \times\frac{n(x',x)}{N_{c}^{2}-1}\Big[ N_{c}{\rm Tr}\Big\{\big(U^{\dagger}(y')-U^{\dagger}(x')\big)\big(U(x)-U(z)\big)\Big\}\notag\\
&\quad\quad\quad -{\rm Tr}\Big\{T^{e}U^{\dagger}(z)\big(U(w)-U(z)\big)T^{e}\big(U^{\dagger}(y')-U^{\dagger}(x')\big)U(z)\Big\}\Big]\notag\\
&\quad\quad -\Big[(x-z)\cdot\Kv(x'-z)\, \Kv(z-w)\cdot\Kv(y'-z) -\frac{1}{2}(x-z)\cdot\Kv(x'-z)\, \Kv(x-w)\cdot\Kv(y'-z)\Big]\notag\\
&\quad\quad \times\frac{n(x',x)}{N_{c}^{2}-1}\Big[ N_{c}{\rm Tr}\Big\{\big(U^{\dagger}(z)-U^{\dagger}(x')\big)\big(U(x)-U(z)\big)\Big\}\notag\\
&\quad\quad\quad -{\rm Tr}\Big\{T^{e}U^{\dagger}(z)\big(U(w)-U(z)\big)T^{e}\big(U^{\dagger}(z)-U^{\dagger}(x')\big)U(z)\Big\}\Big]\bigg\}\notag\\
&\quad +\frac{1}{(2\pi)^{3}}\, \frac{g^{4}}{8\pi^{5}}\, \frac{1}{k^{+}} \int_{w,x',w',z,x}e^{-ik\cdot(w-w')}\, \frac{(x-w)^{2}}{\big[(x-z)^{2}-(x-w)^{2}\big]}\notag\\
&\quad \times \log\left[\frac{(x-z)^{2}}{(x-w)^{2}}\right] \bigg[\Kv(x-z)\cdot\Kv(x'-z)\, \Kv(z-w)\cdot\Kv(w'-z)\notag\\
&\quad\quad -\frac{1}{2}\, \Kv(x-z)\cdot\Kv(x'-z)\, \Kv(x-w)\cdot\Kv(w'-z)\bigg]\notag\\
&\quad \times \frac{n(x',x)}{N_{c}^{2}-1}\Big[ N_{c}{\rm Tr}\Big\{\big(U^{\dagger}(z)-U^{\dagger}(x')\big)\big(U(x)-U(z)\big)\Big\}\notag\\
&\quad\quad -{\rm Tr}\Big\{T^{e}U^{\dagger}(z)\big(U(w)-U(z)\big)T^{e}\big(U^{\dagger}(z)-U^{\dagger}(x')\big)U(z)\Big\}\Big]\, +{\rm c.c.}\, ,
\end{align}
\endgroup


\paragraph{Row III.V.} The symmetrised three-$\rho$ part is
\begingroup\small
\begin{align}\label{eq:S-III.V-3}
&\Big\langle\frac{d^{3}N^{\rm III.V}_{{\rm sym}}}{d^{3}k}\Big\rangle=\frac{1}{(2\pi)^{3}}\, \frac{ig^{4}}{4\pi^{4}}\, \frac{1}{k^{+}} \int_{0}^{\infty}\frac{d\xi}{2\pi} \int_{w,y',z,y,x,x'}e^{-ik\cdot(w-y')}e^{-ik\cdot(z-y')\,\xi}\Bigg[\notag\\
&\quad\quad \mathfrak{K}^{(y')}_{\rm III} \times \frac{N_{c}}{(N_{c}^{2}-4)(N_{c}^{2}-1)}\, \Big(+i\Big)\, \sigma(x',x,y)\Big[\notag\\
&\quad\quad\quad {\rm Tr}\Big\{T^{h}U^{\dagger}(x)\big(U(y')-U(x')\big)D^{h}\big(U^{\dagger}(y)-U^{\dagger}(w)\big)U(x)\Big\}\Big]\notag\\
&\quad\quad + \mathfrak{K}^{(z)}_{\rm III} \times \frac{N_{c}}{(N_{c}^{2}-4)(N_{c}^{2}-1)}\, \Big(+i\Big)\, \sigma(x',x,y)\Big[\notag\\
&\quad\quad\quad {\rm Tr}\Big\{T^{h}U^{\dagger}(x)\big(U(z)-U(x')\big)D^{h}\big(U^{\dagger}(y)-U^{\dagger}(w)\big)U(x)\Big\}\Big]\Bigg] +{\rm c.c.}
\end{align}
\endgroup

and the contact part is
\begingroup\small
\begin{align}\label{eq:S-III.V-2}
&\Big\langle\frac{d^{3}N^{\rm III.V}_{{\rm con}}}{d^{3}k}\Big\rangle=\frac{1}{(2\pi)^{3}}\, \frac{ig^{4}}{4\pi^{4}}\, \frac{1}{k^{+}} \int_{0}^{\infty}\frac{d\xi}{2\pi} \int_{w,y',z,y,x,x'}e^{-ik\cdot(w-y')}e^{-ik\cdot(z-y')\,\xi}\Bigg[\notag\\
&\quad\quad \mathfrak{K}^{(y')}_{\rm III} \times \frac{1}{N_{c}^{2}-1}\bigg[\Big(+\frac{1}{2}i\Big)\delta^{(2)}(x'-x)\, n(x,y)\Big[\notag\\
&\quad\quad\quad\quad {\rm Tr}\Big\{T^{e}U^{\dagger}(x)\big(U(y)-U(w)\big)T^{e}\big(U^{\dagger}(y')-U^{\dagger}(x)\big)U(x)\Big\}\Big]\notag\\
&\quad\quad\quad +\Big(+\frac{1}{2}i\Big)\delta^{(2)}(y-x)\, n(x,x')\Big[ {\rm Tr}\Big\{T^{e}U^{\dagger}(x)\big(U(x)-U(w)\big)T^{e}\big(U^{\dagger}(y')-U^{\dagger}(x')\big)U(x)\Big\}\Big]\notag\\
&\quad\quad\quad +\Big(-\frac{1}{2}i\Big)\delta^{(2)}(x'-y)\, n(y,x)\Big[ {\rm Tr}\Big\{T^{e}U^{\dagger}(x)\big(U(y)-U(w)\big)T^{e}\big(U^{\dagger}(y')-U^{\dagger}(y)\big)U(x)\Big\}\Big]\bigg]\notag\\
&\quad\quad + \mathfrak{K}^{(z)}_{\rm III} \times \frac{1}{N_{c}^{2}-1}\bigg[\Big(+\frac{1}{2}i\Big)\delta^{(2)}(x'-x)\, n(x,y)\Big[\notag\\
&\quad\quad\quad\quad {\rm Tr}\Big\{T^{e}U^{\dagger}(x)\big(U(y)-U(w)\big)T^{e}\big(U^{\dagger}(z)-U^{\dagger}(x)\big)U(x)\Big\}\Big]\notag\\
&\quad\quad\quad +\Big(+\frac{1}{2}i\Big)\delta^{(2)}(y-x)\, n(x,x')\Big[ {\rm Tr}\Big\{T^{e}U^{\dagger}(x)\big(U(x)-U(w)\big)T^{e}\big(U^{\dagger}(z)-U^{\dagger}(x')\big)U(x)\Big\}\Big]\notag\\
&\quad\quad\quad +\Big(-\frac{1}{2}i\Big)\delta^{(2)}(x'-y)\, n(y,x)\Big[ {\rm Tr}\Big\{T^{e}U^{\dagger}(x)\big(U(y)-U(w)\big)T^{e}\big(U^{\dagger}(z)-U^{\dagger}(y)\big)U(x)\Big\}\Big]\bigg]\Bigg]\notag\\
&\quad +{\rm c.c.}
\end{align}
\endgroup


\paragraph{Row III.VI.} The symmetrised three-$\rho$ part is
\begingroup\small
\begin{align}\label{eq:S-III.VI-3}
&\Big\langle\frac{d^{3}N^{\rm III.VI}_{{\rm sym}}}{d^{3}k}\Big\rangle=\frac{1}{(2\pi)^{3}}\, \frac{ig^{4}}{4\pi^{4}}\, \frac{1}{k^{+}} \int_{0}^{\infty}\frac{d\xi}{2\pi} \int_{w,y',z,y,x,x'}e^{-ik\cdot(w-y')}e^{-ik\cdot(z-y')\,\xi}\Bigg[\notag\\
&\quad\quad \mathfrak{K}^{(y')}_{\rm III} \times \frac{N_{c}}{(N_{c}^{2}-4)(N_{c}^{2}-1)}\, \Big(+i\Big)\, \sigma(x',y,x)\Big[\notag\\
&\quad\quad\quad {\rm Tr}\Big\{T^{h}U^{\dagger}(y)\big(U(x)-U(z)\big)D^{h}\big(U^{\dagger}(y')-U^{\dagger}(x')\big)U(y)\Big\}\Big]\notag\\
&\quad\quad + \mathfrak{K}^{(z)}_{\rm III} \times \frac{N_{c}}{(N_{c}^{2}-4)(N_{c}^{2}-1)}\, \Big(+i\Big)\, \sigma(x',y,x)\Big[\notag\\
&\quad\quad\quad {\rm Tr}\Big\{T^{h}U^{\dagger}(y)\big(U(x)-U(z)\big)D^{h}\big(U^{\dagger}(z)-U^{\dagger}(x')\big)U(y)\Big\}\Big]\Bigg] +{\rm c.c.}
\end{align}
\endgroup

and the contact part is
\begingroup\small
\begin{align}\label{eq:S-III.VI-2}
&\Big\langle\frac{d^{3}N^{\rm III.VI}_{{\rm con}}}{d^{3}k}\Big\rangle=\frac{1}{(2\pi)^{3}}\, \frac{ig^{4}}{4\pi^{4}}\, \frac{1}{k^{+}} \int_{0}^{\infty}\frac{d\xi}{2\pi} \int_{w,y',z,y,x,x'}e^{-ik\cdot(w-y')}e^{-ik\cdot(z-y')\,\xi}\Bigg[\notag\\
&\quad\quad \mathfrak{K}^{(y')}_{\rm III} \times \frac{1}{N_{c}^{2}-1}\bigg[\Big(+\frac{1}{2}i\Big)\delta^{(2)}(x'-y)\, n(y,x)\Big[\notag\\
&\quad\quad\quad\quad -{\rm Tr}\Big\{T^{e}U^{\dagger}(y)\big(U(x)-U(z)\big)T^{e}\big(U^{\dagger}(y')-U^{\dagger}(y)\big)U(y)\Big\}\Big]\notag\\
&\quad\quad\quad +\Big(+\frac{1}{2}i\Big)\delta^{(2)}(y-x)\, n(x,x')\Big[ -{\rm Tr}\Big\{T^{e}U^{\dagger}(x)\big(U(x)-U(z)\big)T^{e}\big(U^{\dagger}(y')-U^{\dagger}(x')\big)U(x)\Big\}\Big]\notag\\
&\quad\quad\quad +\Big(-\frac{1}{2}i\Big)\delta^{(2)}(x'-x)\, n(x,y)\Big[ -{\rm Tr}\Big\{T^{e}U^{\dagger}(y)\big(U(x)-U(z)\big)T^{e}\big(U^{\dagger}(y')-U^{\dagger}(x)\big)U(y)\Big\}\Big]\bigg]\notag\\
&\quad\quad + \mathfrak{K}^{(z)}_{\rm III} \times \frac{1}{N_{c}^{2}-1}\bigg[\Big(+\frac{1}{2}i\Big)\delta^{(2)}(x'-y)\, n(y,x)\Big[\notag\\
&\quad\quad\quad\quad -{\rm Tr}\Big\{T^{e}U^{\dagger}(y)\big(U(x)-U(z)\big)T^{e}\big(U^{\dagger}(z)-U^{\dagger}(y)\big)U(y)\Big\}\Big]\notag\\
&\quad\quad\quad +\Big(+\frac{1}{2}i\Big)\delta^{(2)}(y-x)\, n(x,x')\Big[ -{\rm Tr}\Big\{T^{e}U^{\dagger}(x)\big(U(x)-U(z)\big)T^{e}\big(U^{\dagger}(z)-U^{\dagger}(x')\big)U(x)\Big\}\Big]\notag\\
&\quad\quad\quad +\Big(-\frac{1}{2}i\Big)\delta^{(2)}(x'-x)\, n(x,y)\Big[ -{\rm Tr}\Big\{T^{e}U^{\dagger}(y)\big(U(x)-U(z)\big)T^{e}\big(U^{\dagger}(z)-U^{\dagger}(x)\big)U(y)\Big\}\Big]\bigg]\Bigg]\notag\\
&\quad +{\rm c.c.}
\end{align}
\endgroup


\endgroup

\subsection{Physical content of the finite terms}\label{sec:fin_physics}

The finite terms fall into five classes.

\paragraph{(1) A term proportional to the leading order.} The genuine two-$\rho$ term of Row~I.II, Eq.~\eqref{eq:S-I.II-2o}, is $\frac{\alpha_sN_c}{2\pi}(\frac{67}{9}-\frac{\pi^2}{3})$ times the LO spectrum. It is the finite part of the self-energy and vertex corrections of the measured gluon in the wave function. It must not be read as ``the'' NLO constant: by the identity~\eqref{eq:dg_SEid} the self-energy part that carries it is cancelled by the part of Row~I.V in which the target does not resolve the gluon pair, and in the grouping used here the compensating terms are in the genuine two-$\rho$ part of Row~I.V (Section~\ref{sec:lit_constant}).

\paragraph{(2) Recoil corrections to the target evolution.} The plus-prescription terms of Rows~I.III and~I.V (virtual) and of Rows~III.III, III.V and~III.VI (real) have the transverse and colour structure of the JIMWLK term, with the logarithm of the cut-off replaced by a function of the recoil variable. They express the fact that a slow gluon emitted or absorbed by the measured gluon is restricted by the kinematics of the measured gluon: they are the difference between the strict leading-logarithmic evolution of the target and an evolution with a kinematic constraint, Eq.~\eqref{eq:kinconstr}.

\paragraph{(3) Production of two gluons with comparable momenta.} The genuine two-$\rho$ parts of Rows~II.I, II.IV and~III.III are, together, $|\mathcal A_B-\mathcal A_{CA}|^2$ integrated over the second gluon: the square of the amplitude for two final-state gluons, one from the wave function and one from the splitting after the collision, minus its limits $p^+\gg k^+$, $p^+\ll k^+$ and collinear, which are in the BFKL, JIMWLK and DGLAP terms. In the dilute limit this is the two-gluon production cross section of quasi-multi-Regge kinematics minus its multi-Regge limit (Section~\ref{sec:check_dilute}). The remnant and the tensor term of Row~II.IV belong to this class; the latter has the structure of a linearly polarised parent gluon and vanishes after the average over the direction of $\v y'-\v y$.

\paragraph{(4) Terms odd in $\v k$: the symmetrised cubic terms.} They are of first order in the density of the projectile beyond the leading one, are proportional to the C-odd three-point function $\sigma$ of the projectile and to C-odd (odderon-type) traces of three Wilson lines, and are the only source of a difference between the spectra at $\v k$ and $-\v k$. They vanish for $SU(2)$ and for a projectile without a $d^{abc}$ three-point function, and after an average over a C-even target ensemble.

\paragraph{(5) Contact terms.} The commutators of the ordered charges give two-point terms with two charges at one transverse position. They have no counterpart in a calculation with classical charges, are even in $\v k$, and are of the same order as the terms of class~(3); in the dilute limit they are required for the agreement with the known two-gluon production cross section (Section~\ref{sec:check_dilute}). They do not cancel among the rows at the level of the integrand.

\paragraph{Integrability.}\label{sec:fin_summary} With the groupings of Section~\ref{sec:fin_together} the finite terms are integrable at fixed positions of the measured gluon. The statement rests on exact identities, on power counting at all coincidences of the integration points with the colour factors as written, on the pointwise large-distance behaviour in all directions and on Fourier integrals at fixed values of the other coordinates (Appendix~\ref{app:tails} and Section~\ref{sec:lit_finite}); a complete numerical integration of the finite terms over all transverse positions with Wilson lines has not been carried out (Section~\ref{sec:openpoints}).

